\section{The geometric realization and the Galois representations}\label{sec:pf:geo}

\providecommand{\geoES}{\mathrm{ES}}
\providecommand{\geoJ}{\mathrm{J}}
\providecommand{\geoPhiK}{\Phi}
\providecommand{\geoW}{\mathrm{Wt}}
\providecommand{\geoRich}{\mathrm{Rich}}
\providecommand{\geoPar}{\mathrm{Par}}
\providecommand{\geoAbove}{\mathrm{Ab}}
\providecommand{\geotw}{\mathrm{tw}}
\providecommand{\geoHk}{P^{\mathrm{Hk}}}
\providecommand{\geort}{P^{\mathrm{rt}}}
\providecommand{\geounr}{P^{\mathrm{unr}}}
\providecommand{\geoHdg}[3]{H^{#3}\bigl(#1,\Omega^{#2}\bigr)}
\providecommand{\geoGU}{\mathrm{GU}}
\providecommand{\geoUa}{\Uni_a(\A_{E^+})}
\providecommand{\geoHO}{\mathrm{HO}}
\providecommand{\geoCong}{\mathcal{C}}
\providecommand{\geoGood}{\mathcal{G}}
\providecommand{\geoSat}{\mathrm{Sat}}
\providecommand{\geoWD}{\mathrm{WD}}

This section proves three results of Part~I:
\begin{itemize}[nosep]
\item \cref{thm:row0b}: the existence of $\ell$-adic Galois representations, with partial
  local--global compatibility, attached to regular algebraic polarized cuspidal automorphic
  representations of $\GL_a$ over CM fields ($\sfname{Row0b}(\mathrm{recT})$);
\item \cref{thm:geoone}: the geometric realization theorem $\sfname{GeoOne}$ for automorphic
  representations of totally definite unitary groups;
\item \cref{thm:galois:hlf}: the holomorphic Lefschetz formula for finite \'etale correspondences on
  smooth projective complex varieties.
\end{itemize}
It also proves a number of auxiliary statements used in other sections: the ten statements
(S7)--(S13), (S15)--(S17) of $\sfname{Steps}$ (\cref{def:pf:steps2:steps}) which are not proved in
\cref{sec:pf:steps1,sec:pf:steps2}, several facts on $p$-adic Hodge theory and on Hodge
cohomology used there and in \cref{sec:pf:ks1}, identities for Schur polynomials, and Borel's
finiteness of class numbers in the form used.

\paragraph{Structure of the proof.}
The proof of \cref{thm:row0b} is a composition of three results: the reduction
\lref{UpperTheory.Gal.Int.galoisLocalGlobalStrong_of_geometry_lift}, which states that
$\sfname{GeoOne}$ and $\sfname{TLiftK}$ (\cref{def:galois:tliftk}) imply
$\sfname{Row0b}(\mathrm{recT})$; the theta lifting theorem \cref{thm:galois:tlift}; and
\cref{thm:geoone}. The proof of \cref{thm:geoone} has three layers (\cref{sec:pf:geo:geoone}).
A \emph{plan} for an automorphic representation $U$ of the definite unitary group $\Uni_a$ is a
purely group-theoretic object: an auxiliary prime $p$, a PEL datum of Harris--Taylor type for $p$
with its unitary similitude group $G$, a family of Hecke data for $G$, an algebraic Hecke character
$\kappa$, a weight, and a test function whose definite traces detect the Hecke eigensystem of $U$
(the \emph{spectral condition} of \cref{pf:ks1:def-spec}). A \emph{package} is the geometric
counterpart: a smooth projective Hecke tower over $E$ whose correspondences are finite \'etale,
with congruence relations and Hodge occurrences at prescribed places. The Galois representations
are the twists, by the $\ell$-adic characters of $\kappa$, of Galois representations cut out from the
$\ell$-adic cohomology of the tower by the Hecke eigenvalues. Accordingly the proof combines:
\begin{enumerate}[label=(\roman*),nosep]
\item the seventeen step statements $\sfname{Steps}$ (\cref{thm:pf:steps2:main}), which turn a
  package into Galois representations (\cref{pf:geo:packrep,pf:geo:geooneof});
\item the Kuga--Sato package $\sfname{KSAll}$ (\cref{pf:ks2:def:ksall}), which turns a plan into a
  package (\cref{pf:geo:planpack}); it is proved in \cref{pf:ks2:main} from the holomorphic
  Lefschetz formula;
\item the existence of plans $\sfname{PlanEx}$ (\cref{pf:geo:planex}), proved in
  \cref{sec:pf:geo:plans}.
\end{enumerate}

\paragraph{Organisation.}
\Cref{sec:pf:geo:conv} fixes notation. \Cref{sec:pf:geo:row} proves \cref{thm:row0b} from
\cref{thm:geoone}. \Cref{sec:pf:geo:geoone} defines packages and plans and proves
\cref{thm:geoone} from (i)--(iii). \Cref{sec:pf:geo:twist} proves the ten step statements
(S7)--(S13), (S15)--(S17), and \cref{sec:pf:geo:padic} collects the $p$-adic Hodge-theoretic facts
used by the steps. \Cref{sec:pf:geo:hodge} contains the lemmas on Hodge cohomology, idempotents,
Schur polynomials, level groups and class numbers used in the construction of packages.
\Cref{sec:pf:geo:plans} proves the existence of plans. \Cref{sec:pf:geo:hlf} treats the
holomorphic Lefschetz formula.

\paragraph{What is proved here and what is deferred.}
All arguments at the top of the proof are given in full: \cref{thm:row0b,thm:geoone}, the
passage from plans to packages and from packages to Galois representations, the ten step
statements, the assembly of plans from their components (the \emph{master lemma},
\cref{pf:geo:master}), and the reductions in the proof of \cref{thm:galois:hlf}. Several components
of the construction of plans and the central algebraic lemma of the proof of
\cref{thm:galois:hlf} (the \emph{chart-lift theorem}) are stated precisely and cited by their names in
the formal development \cite{LeanSource2026}; the most substantial of them is the spectral
component $\sfname{SpecTop}$, whose architecture is described in \cref{sec:pf:geo:spec}. Classical
theorems which are proved in the formal development (Borel's finiteness of class numbers, the
Hodge--Tate decomposition of de Rham representations, the holomorphic Lefschetz formula of
Atiyah--Bott for maps, the existence of $\ell$-adic characters attached to algebraic Hecke
characters) are cited from the literature.

% ===========================================================================
\subsection{Notation}\label{sec:pf:geo:conv}

Throughout, $E$ is a CM field with maximal totally real subfield $E^+$ and complex conjugation
$c$; $\Omega$ is an algebraic closure of $E$ and $G_E=\Aut_E(\Omega)$; $\ell$ is a prime, $\Qlbar$ the
algebraic closure of $\Ql$ used in the development, and $\iota\colon\Qlbar\to\C$ a field isomorphism.
A \emph{place} of $E$ is a nonzero prime $w$ of $\OO_E$; we write $\Nm w=\#(\OO_E/w)$, $p_w$ for the
rational prime in $w$, and $w^c$ for the conjugate place $c(w)$. A place $w$ is \emph{split} if
some place $u\ne w$ has the same restriction to $\OO_{E^+}$ (then $u=w^c$), and two places are
\emph{apart} if their restrictions to $\OO_{E^+}$ differ. For $x\in E$ with $x^2=-d$, $d>0$, two
places $w,w'$ \emph{induce the same prime of $\Q(x)$}, written $w\sim_xw'$, if an element of
$\OO_E\cap\Q(x)$ lies in $w$ if and only if it lies in $w'$. This is an equivalence relation, and
$w\sim_xw'$ implies $p_w=p_{w'}$ and $w^c\sim_xw'^c$. A place $w$ is \emph{$d$-usable} if it is split
and $-d$ is a square in $\Q_{p_w}$. The field condition $\sfname{FieldOK}(E,\ell,d)$ of
\cref{def:galois:geoone} means: $d>0$ and $-d$ is a square in $E$; $-d$ is a square in $\Ql$;
$[E^+:\Q]$ is even; and every non-split place $w$ satisfies $\ell\notin w$ and is unramified over
$E^+$. For an embedding $\tau\colon E\to\C$ and a place $v$ with $\ell\in v$ we say that $\tau$
\emph{lies over $v$} if $|\iota^{-1}\tau(y)|<1$ for all $y\in v$.

For $a\ge1$, $\geoUa=\{g\in\GL_a(\A_E):{}^t\bar gg=1\}$ is the group of adelic points of the
totally definite unitary group of \cref{sec:galois:arch}. We use the following notions for a
$\C$-linear space $U$ of functions on $\geoUa$ (\cref{def:galois:geoone},
\cref{def:pf:steps2:lac}): $U$ is \emph{automorphic} if it is a nonzero irreducible space of
automorphic forms stable under right translations; for a split place $w$ and a multiset $B$ of $a$
nonzero complex numbers, $B$ is a \emph{Satake multiset of $U$ at $w$}, written
$\geoSat(U,w,B)$, if some smooth irreducible representation of $\GL_a(E_w)$ which is spherical with
Satake parameters $B$ embeds into $U$ compatibly with the local factor at $w$; a dominant
$\lambda\in\Z^a$ is a \emph{weight of $U$ at $\tau$} if the algebraic representation of highest weight
$\lambda$ occurs in $U$ along the archimedean factor at $\tau$; $J(\lambda)=\{k-\lambda_k:0\le
k<a\}$; and the conditions ``unramified at the non-split places outside $T$'', ``supercuspidal
type of $\ell$-power order at $w_0$'', ``spherical at the split places apart from $w_0$ and from $S$''
and ``lacunary weight at $v'$'' are those of \cref{def:galois:geoone}.

On the Galois side we use the notions of \cref{def:pf:steps2:es,def:pf:steps2:crys,def:pf:steps2:lac}.
For $\rho\colon G_E\to\GL_{dd}(\Qlbar)$, a place $w$ and multisets $R\subset\C$:
$\geoES_w(\rho,R)$ is the Eichler--Shimura condition (for every place $A$ of $\Omega$ over $w$, $\rho$ is
trivial on $I_A$ and $(\iota^{-1}\geort_R)(\rho(\sigma^{-1}))=0$ for every arithmetic Frobenius
element $\sigma$ for $(A,\Nm w)$, where $\geort_R=\prod_{r\in R}(X-r)$); $\rho$ is
\emph{crystalline at $w$}; $\geoJ_w(\rho,Q)$ is the condition that the jumps of $\rho$ at $w$
contain the integers $j$ with $Q(\iota\circ\tau,j)$; $\geoPhiK_w(\rho,R)$ means that $\rho$ is
crystalline at $w$ and that the characteristic polynomial of some element of degree $-1$ of the
Weil group on the Weil--Deligne representation of $\rho$ at $w$ divides a power of $\geort_R$;
$\geoW_w(U,\rho)$ means that for all weights $\lambda$ of $U$ at embeddings over $w$ the set
$J(\lambda)$ is contained in a jump multiset of $\rho$; and $\geoRich_a(\rho)$ means that at some
place above $\ell$ and some embedding the jumps of $\rho$ contain a lacunary set of $a$ integers. All
these conditions quantify over the auxiliary data (period rings, places $A$, bridges) as explained in
\cref{def:pf:steps2:crys}. Finally, for $q\in\N$ and a multiset $B$ we put
$q^{(a-1)/2}B=\{q^{(a-1)/2}b:b\in B\}$, and for $B$ a multiset of $a$ nonzero complex numbers we write
$\geoPhiK^{\mathrm{sat}}_w(\rho;a,B)$ for the condition $\geoPhiK_w$ in which $\geort_R$ is replaced by
the characteristic polynomial of the same Weil group element on the unramified Weil--Deligne
representation with Frobenius eigenvalues $B^{-1}$ twisted by $|\cdot|^{(1-a)/2}$.

\begin{definition}[Basic representations]\label{pf:geo:vbase}
Let $E,\ell,\iota$ be as above, $a,d\in\N$, $U$ a space of functions on $\geoUa$ and
$\rho\colon G_E\to\GL_{dd}(\Qlbar)$. Then $\rho$ is \emph{basic} for $(\iota,d,U)$ if: $\rho$ is continuous
and semisimple; there is a finite set $X$ of places such that $\geoES_w(\rho,(\Nm w)^{(a-1)/2}B)$ holds
for every $d$-usable place $w\notin X$ with $\ell\notin w$ and every Satake multiset $B$ of $U$ at $w$;
and $\geoRich_a(\rho)$ holds.
\formal{UpperTheory.Gal.Int.Div.VBase}
\end{definition}

This is the notion of \cref{def:galois:geoone}, conditions (1) and (2).

% ===========================================================================
\subsection{The Galois representations}\label{sec:pf:geo:row}

\begin{theorem}[\cref{thm:row0b}]\label{pf:geo:row0b}
The proposition $\sfname{Row0b}(\mathrm{recT})$ of \cref{def:row0b} holds.
\formallabel{UpperTheory.Gal.Int.galoisLocalGlobalStrong_recParameter}
\end{theorem}

The recognition notion $\mathrm{recT}$ is the predicate ``is a local parameter map'' of
\cref{def:rectexists}.

\begin{proposition}\label{pf:geo:ofthetalift}
$\sfname{TLiftK}$ implies $\sfname{Row0b}(\mathrm{recT})$.
\formallabel{UpperTheory.Gal.Int.galoisRepresentations_of_thetaLift}
\end{proposition}

\begin{proof}
By \lref{UpperTheory.Gal.Int.galoisLocalGlobalStrong_of_geometry_lift}, the conjunction of
$\sfname{GeoOne}$ and $\sfname{TLiftK}$ implies $\sfname{Row0b}(\mathrm{recT})$; this is
\cref{thm:galois:reduction}, whose proof (descent, congruences, limits of pseudo-characters, theta
lifting from odd to even rank) is given in \cref{sec:pf:galred}. $\sfname{GeoOne}$ holds by \cref{pf:geo:geoone}.
\end{proof}

\begin{proof}[Proof of \cref{pf:geo:row0b}]
$\sfname{TLiftK}$ holds by \cref{thm:galois:tlift}, i.e.\ by \lref{UpperTheory.TP.K.tLiftK},
the existence of theta lifts between definite unitary groups with the archimedean types and
local correspondences of \cref{def:galois:tliftk}. Apply \cref{pf:geo:ofthetalift}.
\end{proof}

\begin{remark}\label{pf:geo:rem:rowinputs}
The proof of \cref{thm:row0b} thus rests on exactly three results: \cref{thm:galois:reduction},
\cref{thm:galois:tlift} and \cref{thm:geoone}. In particular neither the local Langlands
correspondence nor any published existence theorem for Galois representations is used; the local
parameter maps enter only through the laws of \cref{def:rectexists} (\cref{rem:galois:suffices}).
\end{remark}

% ===========================================================================
\subsection{Packages, plans and the geometric realization}\label{sec:pf:geo:geoone}

\subsubsection{The statement}

We recall \cref{thm:geoone} in the form in which it is proved; the hypotheses are those of
\cref{def:galois:geoone}.

\begin{theorem}[\cref{thm:geoone}]\label{pf:geo:geoone}
Let $E$ be a CM field, $\ell$ a prime, $\iota\colon\Qlbar\to\C$ a field isomorphism, $a$ an odd natural
number, $d\in\N$ with $\sfname{FieldOK}(E,\ell,d)$ and $x\in E$ with $x^2=-d$. Let $T$ be a finite set of
non-split places, $w_0,v'$ places containing $\ell$ which are apart and satisfy $w_0\sim_xv'$, $S$ a
finite set of places and $U$ a space of functions on $\geoUa$ which is automorphic and unramified at
the non-split places outside $T$, contains a supercuspidal type of $\ell$-power order at $w_0$, has
Satake multisets at all split places apart from $w_0$ and from all places of $S$, has a Satake
multiset at $v'$, and has a lacunary weight at $v'$. Then:
\begin{enumerate}[label=\textup{(G\arabic*)},nosep]
\item there is a basic representation $\rho\colon G_E\to\GL_{dd}(\Qlbar)$ for $(\iota,d,U)$;
\item for every $d$-usable place $w$ with $\ell\notin w$ at which $U$ has a Satake multiset there is a
  basic $\rho$ with $\geoES_w(\rho,(\Nm w)^{(a-1)/2}B)$ for every Satake multiset $B$ of $U$ at $w$;
\item for every place $w$ with $\ell\in w$, $w\sim_xv'$, $w$ apart from $w_0$ and from $v'$, at which
  $U$ has a Satake multiset: there is a basic $\rho$ which is crystalline at $w$ and satisfies
  $\geoW_w(U,\rho)$; and there is a basic $\rho$ with $\geoPhiK^{\mathrm{sat}}_w(\rho;a,B)$ for every
  Satake multiset $B$ of $U$ at $w$.
\end{enumerate}
\formallabel{UpperTheory.Gal.Int.Geo.geoOne_holds}\formal{UpperTheory.Gal.Int.Geo.GeoOne}
\end{theorem}

The three conclusions are separate existence statements: for each place a basic representation
with the local property at that place is provided, but not one representation with all properties.
The dimension $dd$ is not prescribed; the passage to $a$-dimensional representations is part of the
proof of \cref{thm:galois:reduction}.

\subsubsection{Packages}

\begin{definition}[Packages]\label{pf:geo:pack}
Let $E,\ell,\iota$ be as above, $a,d\in\N$ and $U$ a space of functions on $\geoUa$. A \emph{package}
$P$ for $(\iota,d,U)$ consists of a nice Hecke tower $H=(Z,(Y_g,s_g,t_g)_{g\in\Gamma})$ over $E$
(\cref{def:pf:steps2:tower}), a map $\theta\colon\Gamma\to\Qlbar$, a Hecke character $\kappa$ of $E$ and
a map $m\colon\Hom(E,\C)\to\Z$ such that $\kappa$ is algebraic of infinity type $m$ (that is,
$\kappa_\infty(k)=\prod_\phi\phi(k)^{m(\phi)}$ for every totally positive $k\in E^\times$), and a finite set
$X$ of places such that $\sfname{CongFor}(w,B)$ holds for every $d$-usable $w\notin X$ with $\ell\notin
w$ and every Satake multiset $B$ of $U$ at $w$. Here $\sfname{CongFor}(w,B)$ means: $\kappa$ is
unramified at $w$, and there is $\vartheta\colon\N\to\Qlbar$ such that $(H,\theta)$ satisfies the
congruence relation at $w$ with coefficients $\vartheta$ in degree $a$
($\geoCong(H,\theta,w,a,\vartheta)$, \cref{pf:ks2:def:goodatw}) and
$\kappa(w)^{-1}(\Nm w)^{(a-1)/2}B$ is an $\iota$-root multiset for $(w,a,\vartheta)$
(\cref{def:pf:steps2:es}). For a place $v$:
\begin{itemize}[nosep]
\item $\mathrm{HTFor}(P,v)$ means: $\kappa$ is unramified at $v$, $(H,\theta)$ has good reduction at
  $v$ ($\geoGood(H,\theta,v)$, \cref{pf:ks2:def:goodatw}), and for every embedding $\tau\colon E\to\C$
  over $v$ there is $\sigma\in\Aut(\C)$ such that for every weight $\lambda$ of $U$ at $\tau$ and every
  $j\in J(\lambda)$ the Hodge predicate $\mathrm{Hdg}_{H,\theta}$ for the isomorphism $\sigma\circ\iota$ holds
  at $(\sigma\circ\tau,j+m(\tau))$: $j+m(\tau)\ge0$ and the eigensystem $\sigma\circ\iota\circ\theta$ occurs in
  some $\geoHdg{Z\otimes_{E,\sigma\tau}\C}{j+m(\tau)}{b}$ (\cref{def:pf:steps2:crys});
\item $\mathrm{CongAll}(P,w)$ means: $\sfname{CongFor}(w,B)$ for every Satake multiset $B$ of $U$ at $w$.
\end{itemize}
$\sfname{PackEx}$ is the proposition obtained from \cref{pf:geo:geoone} by deleting the hypothesis of a
lacunary weight and replacing the conclusions by: there is a package $P$ with $\mathrm{HTFor}(P,v')$;
for every $w$ as in (G2) there is a package $P$ with $\mathrm{HTFor}(P,v')$ and $\mathrm{CongAll}(P,w)$;
and for every $w$ as in (G3) there are packages $P,P'$ with $\mathrm{HTFor}(P,v')$, $\mathrm{HTFor}(P,w)$,
$\mathrm{HTFor}(P',v')$ and $\mathrm{CongAll}(P',w)$.
\formal{UpperTheory.Gal.Int.Geo.Pack}\formal{UpperTheory.Gal.Int.Geo.CongFor}\formal{UpperTheory.Gal.Int.Geo.Pack.HTFor}\formal{UpperTheory.Gal.Int.Geo.Pack.CongAll}\formal{UpperTheory.Gal.Int.Geo.PackEx}
\end{definition}

\begin{proposition}[Packages give representations]\label{pf:geo:packrep}
Assume $\sfname{Steps}$. Let $P=(H,\theta,\kappa,m,X)$ be a package for $(\iota,d,U)$ and $v$ a place with
$\ell\in v$ such that $U$ has a lacunary weight at $v$ and $\mathrm{HTFor}(P,v)$ holds. Then there is a
basic representation $\rho\colon G_E\to\GL_{dd}(\Qlbar)$ for $(\iota,d,U)$ such that:
\begin{enumerate}[label=(\alph*),nosep]
\item $\geoES_w(\rho,(\Nm w)^{(a-1)/2}B)$ for every place $w$ with $\ell\notin w$ and $\mathrm{CongAll}(P,w)$
  and every Satake multiset $B$ of $U$ at $w$;
\item $\rho$ is crystalline at $w$ and $\geoW_w(U,\rho)$ holds, for every place $w$ with $\ell\in w$ and
  $\mathrm{HTFor}(P,w)$;
\item $\geoPhiK^{\mathrm{sat}}_w(\rho;a,B)$ for every place $w$ with $\ell\in w$ and $\mathrm{CongAll}(P,w)$ and
  every Satake multiset $B$ of $U$ at $w$.
\end{enumerate}
\formallabel{UpperTheory.Gal.Int.Geo.Pack.rep}
\end{proposition}

\begin{proof}
By (S1) there are $dd$ and a continuous semisimple $\rho_0\colon G_E\to\GL_{dd}(\Qlbar)$ attached to
$(H,\theta)$ (\cref{def:pf:steps2:piece}). By (S15) there is a family $\mathcal L$ of characters
$G_E\to\Qlbar^\times$ such that $\chi_\kappa:=\mathcal L(\kappa,m)$ is continuous and, at every place $w$
where $\kappa$ is unramified, the $1\times1$ representation $\chi_\kappa$ has Frobenius eigenvalue
$\kappa(w)$ if $\ell\notin w$, and is crystalline with jump $-m(\iota\circ\tau)$ at every $\tau$ and with
Weil--Deligne Frobenius eigenvalue $\kappa(w)$ if $\ell\in w$; in the notation of \cref{sec:pf:geo:twist}
this is $\geoPar_w(\chi_\kappa;\{\kappa(w)\},\tau\mapsto\{-m(\iota\circ\tau)\})$. Put
\[
  \rho=\geotw(\rho_0,\chi_\kappa),\qquad \rho(g)=\chi_\kappa(g)\rho_0(g).
\]
By (S17) $\rho$ is continuous and by (S16) it is semisimple. We use the twisting statements
(S7)--(S10) with the character $\chi_\kappa$, the constant $c=\kappa(w)$ and the infinity type
$-m$; their hypothesis is the property of $\chi_\kappa$ at $w$ just recalled.

\emph{Places away from $\ell$.} Let $\ell\notin w$ and $B$ be such that $\sfname{CongFor}(w,B)$ holds,
with $\vartheta$ as in its definition. By (S6) the congruence datum is realized by an onto cover, i.e.\
$\geoCong(H,\theta,w,a,\vartheta)$ in the sense of \cref{def:pf:steps2:cong}. Since
$R=\kappa(w)^{-1}(\Nm w)^{(a-1)/2}B$ is an $\iota$-root multiset for $(w,a,\vartheta)$, (S2) gives
$\geoES_w(\rho_0,R)$, and (S7) gives $\geoES_w(\rho,\kappa(w)R)=\geoES_w(\rho,(\Nm w)^{(a-1)/2}B)$.
Applied to the places outside the finite set $X$ of the package, this gives the almost-everywhere
clause of \cref{pf:geo:vbase}; applied to a place with $\mathrm{CongAll}(P,w)$ it gives (a).

\emph{Places above $\ell$ with good reduction.} Let $\ell\in w$ and assume $\mathrm{HTFor}(P,w)$. By (S6),
$\geoGood(H,\theta,w)$ yields an onto good cover over $\OO_{(w)}$, so $\rho_0$ is crystalline at $w$ by
(S3) and $\rho$ is crystalline at $w$ by (S8). For every embedding $\tau\colon E\to\C$ choose
$\sigma_\tau\in\Aut(\C)$ as in $\mathrm{HTFor}(P,w)$ if $\tau$ lies over $w$, and $\sigma_\tau=\mathrm{id}$
otherwise. Statement (S4) holds for every isomorphism $\Qlbar\to\C$; applying it, for each embedding
$\tau_0\colon E\to\Qlbar$ induced by a bridge, with the isomorphism $\sigma_\tau\circ\iota$ where
$\tau=\iota\circ\tau_0$, we obtain $\geoJ_w(\rho_0,Q)$ for the predicate
\[
  Q(\tau,j)\iff j\ge0\text{ and }\sigma_\tau\iota\theta\text{ occurs in some }
  \geoHdg{Z\otimes_{E,\sigma_\tau\tau}\C}{j}{b}.
\]
By (S9) with the infinity type $-m$, $\geoJ_w(\rho,Q')$ holds with $Q'(\tau,j)=Q(\tau,j+m(\tau))$. By the
choice of $\sigma_\tau$, $Q'(\tau,j)$ holds for every $\tau$ over $w$, every weight $\lambda$ of $U$ at $\tau$
and every $j\in J(\lambda)$, so (S13) gives $\geoW_w(U,\rho)$. This proves (b). Applied at $w=v$
together with the lacunary weight of $U$ at $v$, (S14) gives $\geoRich_a(\rho)$; so $\rho$ is basic.

\emph{Crystalline Frobenius.} Let $\ell\in w$, $\mathrm{CongAll}(P,w)$, and $B$ a Satake multiset of $U$
at $w$; let $\vartheta$ be as in $\sfname{CongFor}(w,B)$. As before (S6) and (S5) give
$\geoPhiK_w(\rho_0,\kappa(w)^{-1}(\Nm w)^{(a-1)/2}B)$, (S10) gives $\geoPhiK_w(\rho,(\Nm w)^{(a-1)/2}B)$,
and since $B$ has $a$ elements by (S12), (S11) gives $\geoPhiK^{\mathrm{sat}}_w(\rho;a,B)$. This is (c).
\end{proof}

\begin{theorem}\label{pf:geo:geooneof}
$\sfname{Steps}$ and $\sfname{PackEx}$ imply $\sfname{GeoOne}$.
\formallabel{UpperTheory.Gal.Int.Geo.geoOne_of}
\end{theorem}

\begin{proof}
Let the data and hypotheses of \cref{pf:geo:geoone} be given. The hypotheses of $\sfname{PackEx}$ are
among them, so $\sfname{PackEx}$ provides the packages listed in \cref{pf:geo:pack}; each satisfies
$\mathrm{HTFor}(\cdot,v')$. Apply \cref{pf:geo:packrep} with $v=v'$, using the lacunary weight of $U$
at $v'$. The package with $\mathrm{HTFor}(P,v')$ gives a basic representation, which is (G1). For $w$ as
in (G2), the package with $\mathrm{CongAll}(P,w)$ gives a basic $\rho$ satisfying (a) at $w$, which is
(G2). For $w$ as in (G3), the package $P$ with $\mathrm{HTFor}(P,w)$ gives, by (b), a basic $\rho$
crystalline at $w$ with $\geoW_w(U,\rho)$, and the package $P'$ with $\mathrm{CongAll}(P',w)$ gives, by
(c), a basic $\rho'$ with $\geoPhiK^{\mathrm{sat}}_w(\rho';a,B)$ for all Satake multisets $B$ at $w$.
\end{proof}

\subsubsection{Plans}

The definition of plans uses the objects of the Kuga--Sato package: data, unitary similitude
groups $G$ and links $L$ of Harris--Taylor type for a prime $p$ with the standing conditions
$\mathrm{Sc}(D,G)$, levels $N$ (integers $N\ge3$ prime to $p$), the principal congruence subgroups
$K_N\subseteq G(\A_f^p)$, Hecke data $g=(K_g,h_g,M_g)$ of level $N$ and their Hecke functions $f_g$, and
the group-theoretic conditions $\mathcal U(g,\theta,\mathfrak w)$ (\emph{unramified at $\mathfrak w$}) and
$\mathcal H(g,\theta,\mathfrak w,\vartheta)$ (\emph{contains the operators $T_i$ at $\mathfrak w$ with
eigenvalues $\vartheta$}); see \cref{pf:ks2:def:datum,pf:ks2:def:heckedata,pf:ks2:def:goodgrp}. It also uses
the spectral condition $\mathcal S$ and Hodge exponent vectors of \cref{pf:ks1:def-spec,pf:ks1:def-hexp},
and the map $\varepsilon_D$ of \cref{pf:ks1:def-hexp} ($\varepsilon_D(\tau)=1$ if and only if $\tau(\omega)$ is the
complex conjugate of the distinguished root $\rho_c$ of the datum).

\begin{definition}[Plans]\label{pf:geo:plan}
Let $E,\ell,\iota,a,d,U$ be as in \cref{pf:geo:pack}. A \emph{plan} $P$ for $(\iota,d,U)$ consists of: a
prime $p$; a datum $D$, a group $G$ and a link $L$ for $p$ with $\mathrm{Sc}(D,G)$; a ring isomorphism
$e\colon F_D\to E$ with $n_D=a$; a level $N$; a family $(g_j)_{j\in J}$ of Hecke data of level $N$ with
values $\theta\colon J\to\Qlbar$; a Hecke character $\kappa$ of $E$ which is algebraic of infinity type
$m\colon\Hom(E,\C)\to\Z$; a weight $\lambda\colon\Hom(F_D,\C)\to\N^{a}$ with each $\lambda(\tau)$
non-increasing; a function $f$ on $G(\A_f^p)$ which lies in the Hecke algebra of $K_N$ and is supported
on elements $y$ with $y^{-1}$ integral; such that the spectral condition
$\mathcal S(L,\lambda,K_N,(f_{g_j})_j,\iota\circ\theta,f)$ holds; and a finite set $X$ of places such that
$\sfname{CongP}(w,B)$ holds for every $d$-usable $w\notin X$ with $\ell\notin w$ and every Satake multiset
$B$ of $U$ at $w$. Here $\sfname{CongP}(w,B)$ means: $\kappa$ is unramified at $w$ and there are a
maximal ideal $\mathfrak w$ of $\OO_{F_D}$ corresponding to $w$ under $e$ and $\vartheta\colon\N\to\Qlbar$ with
$\mathcal H(g,\theta,\mathfrak w,\vartheta)$ and such that $\kappa(w)^{-1}(\Nm w)^{(a-1)/2}B$ is an
$\iota$-root multiset for $(w,a,\vartheta)$. For a place $v$:
\begin{itemize}[nosep]
\item $\mathrm{HTFor}(P,v)$ means: $\kappa$ is unramified at $v$; $\mathcal U(g,\theta,\mathfrak v)$ for a
  maximal ideal $\mathfrak v$ corresponding to $v$; and (\emph{exponent condition}) for every $\tau\colon
  E\to\C$ over $v$ there is $\sigma\in\Aut(\C)$ with $\sigma\tau e=\tau_i$ for all the distinguished embeddings
  $\tau_i$ of $D$, such that for every weight $\lambda_U$ of $U$ at $\tau$ and every $j\in J(\lambda_U)$
  there are an index $i_0$, a Hodge exponent vector $h$ for $(\tau_{i_0},\varepsilon_D,\sigma\lambda)$ and an
  index $i$ with $h_i=j+m(\tau)$; here $(\sigma\lambda)(\tau')=\lambda(\sigma^{-1}\tau')$;
\item $\mathrm{CongAll}(P,w)$ means: $\sfname{CongP}(w,B)$ for every Satake multiset $B$ of $U$ at $w$.
\end{itemize}
$\sfname{PlanEx}$ is $\sfname{PackEx}$ with packages replaced by plans.
\formal{UpperTheory.Gal.Int.Geo.Plan}\formal{UpperTheory.Gal.Int.Geo.CongP}\formal{UpperTheory.Gal.Int.Geo.Plan.HTFor}\formal{UpperTheory.Gal.Int.Geo.Plan.ExpAt}\formal{UpperTheory.Gal.Int.Geo.Plan.CongAll}\formal{UpperTheory.Gal.Int.Geo.PlanEx}
\end{definition}

A plan does not mention schemes. Its geometric realization is the content of the next result.

\begin{proposition}[Plans give packages]\label{pf:geo:planpack}
Assume $\sfname{KSAll}$. For every plan $P$ for $(\iota,d,U)$ there is a package $Q$ for $(\iota,d,U)$
such that $\mathrm{HTFor}(P,v)$ implies $\mathrm{HTFor}(Q,v)$ and $\mathrm{CongAll}(P,w)$ implies
$\mathrm{CongAll}(Q,w)$, for all places $v,w$.
\formallabel{UpperTheory.Gal.Int.Geo.Plan.pack}
\end{proposition}

\begin{proof}
Write $P=(p,D,G,L,e,N,(g_j),\theta,\kappa,m,\lambda,f,X)$. Since $\mathrm{Sc}(D,G)$ includes $a_D=1$,
$\sfname{StdEx}$ gives a characteristic-zero base $g_b$ over $E$ which is standard for $e$. By
$\sfname{KSHodge}$ applied to $(D,G,L)$ and $\lambda$ there is $k\in\N$ with the Hodge occurrence property
for Kuga--Sato towers with exponent $k$, and by $\sfname{KSEx}$ there is a Kuga--Sato tower $T$ for
$(L,g_b,k)$ over $E$. Let $H=T.\mathrm{tower}(N,g)$ be the Hecke tower of the family $g$
(\cref{pf:ks2:def:heckedata}); it is nice by $\sfname{KSNice}$. Put $Q=(H,\theta,\kappa,m,X)$.

If $\sfname{CongP}(w,B)$ holds with $\mathfrak w,\vartheta$, then $\sfname{KSCong}$ (applied with the
standard base $g_b$, the tower $T$ and $\mathfrak w$ corresponding to $w$) turns
$\mathcal H(g,\theta,\mathfrak w,\vartheta)$ into $\geoCong(H,\theta,w,n_D,\vartheta)$, and $n_D=a$; together with
the unramifiedness of $\kappa$ and the root condition, which are the same in both definitions, this is
$\sfname{CongFor}(w,B)$. Applied to the places outside $X$ this shows that $Q$ is a package, and applied
to all Satake multisets at $w$ it shows that $\mathrm{CongAll}(P,w)$ implies $\mathrm{CongAll}(Q,w)$.

Assume $\mathrm{HTFor}(P,v)$. Unramifiedness of $\kappa$ at $v$ is common to both definitions, and
$\sfname{KSGood}$ turns $\mathcal U(g,\theta,\mathfrak v)$ into $\geoGood(H,\theta,v)$. Let $\tau\colon E\to\C$ lie over
$v$ and let $\sigma$ be as in the exponent condition. Let $\lambda_U$ be a weight of $U$ at $\tau$ and
$j\in J(\lambda_U)$; take $i_0,h,i$ with $h$ a Hodge exponent vector for $(\tau_{i_0},\varepsilon_D,\sigma\lambda)$ and
$h_i=j+m(\tau)$. The ring map $\sigma\circ\tau\colon E\to\C$ satisfies $(\sigma\tau)(\rho)=\rho_c$ and
$(\sigma\tau)\circ t_i=\tau_i$ for the components $\rho,t_i$ of $g_b$, because $g_b$ is standard for $e$
($\rho=e(\omega)$, $t_i=e|_{\OO_{F}}$) and $\sigma\tau e=\tau_i$, $\tau_i(\omega)=\rho_c$. Hence $\sfname{KSHodge}$, applied to
$T$, $N$, $g$, the values $\iota\circ\theta$, the function $f$ with its spectral condition, the ring map
$\sigma\tau$, the automorphism $\sigma$, the index $i_0$ and the vector $h$, shows that the eigensystem
$\sigma\circ\iota\circ\theta$ occurs in the Hodge cohomology of the base change of $H$ along $\sigma\tau$ in
degree $h_i=j+m(\tau)$. This is the Hodge predicate required in $\mathrm{HTFor}(Q,v)$ (note $h_i\ge0$).
\end{proof}

\begin{corollary}\label{pf:geo:packexof}
$\sfname{KSAll}$ and $\sfname{PlanEx}$ imply $\sfname{PackEx}$.
\formallabel{UpperTheory.Gal.Int.Geo.packEx_of}
\end{corollary}

\begin{proof}
Given the data of $\sfname{PackEx}$, take the plans provided by $\sfname{PlanEx}$ and replace each by the
package of \cref{pf:geo:planpack}; the properties $\mathrm{HTFor}$ and $\mathrm{CongAll}$ are preserved.
\end{proof}

\subsubsection{Proof of the geometric realization theorem}

\begin{theorem}\label{pf:geo:parts}
$\sfname{Steps}$, $\sfname{KSAll}$ and $\sfname{PlanEx}$ imply $\sfname{GeoOne}$. Consequently the holomorphic
Lefschetz formula $\sfname{HLF}$ (the statement of \cref{thm:galois:hlf}) implies $\sfname{GeoOne}$.
\formallabel{UpperTheory.Gal.Int.Geo.geoOne_of_parts}\formallabel{UpperTheory.Gal.Int.Geo.geometric_realization_of_weight_one}
\end{theorem}

\begin{proof}
The first statement is \cref{pf:geo:geooneof} combined with \cref{pf:geo:packexof}. For the second,
$\sfname{Steps}$ holds by \cref{thm:pf:steps2:main}; $\sfname{HLF}$ implies $\sfname{KSAll}$ by
\cref{pf:ks2:main} (\lref{UpperTheory.Gal.Int.Geo.kodaira_spencer_of_weight_one}: the holomorphic
Lefschetz formula for finite \'etale correspondences implies the Kuga--Sato package); and
$\sfname{PlanEx}$ holds by \cref{pf:geo:planex}.
\end{proof}

\begin{proof}[Proof of \cref{pf:geo:geoone}]
$\sfname{HLF}$ holds by \cref{pf:geo:hlf}; apply \cref{pf:geo:parts}.
\end{proof}

In this proof the holomorphic Lefschetz formula is used only through $\sfname{KSHodge}$; the statements involved do not concern
modular forms of weight one, although the formal names of $\sfname{HLF}$ and of the implication $\sfname{HLF}\Rightarrow\sfname{GeoOne}$
refer to ``weight one''.

% ===========================================================================
\subsection{Twists and normalisations: the statements (S7)--(S13), (S15)--(S17)}\label{sec:pf:geo:twist}

In this subsection $E$ is a number field (a CM field in (S12), (S13)), $\Omega$, $G_E$, $\ell$, $\iota$ are as in
\cref{sec:pf:geo:conv}, and $w$ is a place of $E$. For a homomorphism $\rho\colon\Gamma\to\GL_{dd}(k)$ of a
group to the invertible matrices over a field and a character $\chi\colon\Gamma\to k^\times$, the
\emph{twist} is $\geotw(\rho,\chi)(g)=\chi(g)\rho(g)$, and $\chi^{(1)}$ denotes $\chi$ viewed as a
$1\times1$ matrix representation. For a multiset $S$ of complex numbers and $n\in\Z$ put
\[
  \geounr_S(n)=\prod_{z\in S}\bigl(X-z^{-n}\bigr)\in\C[X],\qquad \geort_S=\prod_{z\in S}(X-z),
\]
so that $\geounr_S(-1)=\geort_S$. The Weil--Deligne representation attached to a potentially
semistable representation of a decomposition group (via a bridge and $\iota$, \cref{def:galois:awayb}(d)) is
written $\geoWD(\cdot)$, and $\mathrm{cp}_g(x)$ is the characteristic polynomial of $g$ on a
Weil--Deligne representation $x$.

\begin{definition}[Local parameters of a representation]\label{pf:geo:par}
Let $r\colon G_E\to\GL_n(\Qlbar)$, $S$ a multiset of complex numbers and $J$ a map from the embeddings
$E\to\Qlbar$ to multisets of integers. Then $\geoPar_w(r;S,J)$ means:
\begin{enumerate}[label=(\alph*),nosep]
\item if $\ell\notin w$: for every place $A$ of $\Omega$ over $w$, $r$ is trivial on $I_A$, and for every
  arithmetic Frobenius element $\sigma$ for $(A,\Nm w)$ the characteristic polynomial of $r(\sigma^{-1})$,
  transported by $\iota$, is $\geort_S$;
\item for every tied period package, every place $A$ over $w$ and every bridge inducing $A$
  (\cref{def:pf:steps2:crys}): $r|_{D_A}$ is crystalline and de Rham, $J(\tau)$ is a jump multiset of
  $r|_{D_A}$ at every embedding $\tau$ induced by the bridge, and
  $\mathrm{cp}_g(\geoWD(r|_{D_A}))=\geounr_S(\deg g)$ for every $g\in W_A$.
\end{enumerate}
\formal{UpperTheory.Gal.TP.Sub.Cl}\formal{UpperTheory.Gal.TP.Sub.Above}\formal{UpperTheory.V.Af.FrobEig}
\end{definition}

Condition (b) is empty if $\ell\notin w$, since a bridge inducing $A$ forces $A$ to lie over $\ell$.

\subsubsection{Local twisting lemmas}

\begin{lemma}\label{pf:geo:twloc}
Let $A$ be a place of $\Omega$ and fix a tied period package and a bridge for $D_A$. Let
$r\colon D_A\to\GL_n(\Qlbar)$ and $c\colon D_A\to\Qlbar^\times$ be homomorphisms.
\begin{enumerate}[label=(\alph*),nosep]
\item If $r$ and $c^{(1)}$ are de Rham (resp.\ crystalline), so is $\geotw(r,c)$.
\item If $r$ and $c$ are continuous, $J$ is a jump multiset of $r$ at $\tau$ and $\{m\}$ is a jump multiset
  of $c^{(1)}$ at $\tau$, then $J+m=\{j+m:j\in J\}$ is a jump multiset of $\geotw(r,c)$ at $\tau$.
\end{enumerate}
\formallabel{UpperTheory.Gal.TP.Sub.twLoc_holds}\formallabel{UpperTheory.Gal.TP.Sub.twJumps_holds}
\end{lemma}

\begin{proof}
Both statements are instances of \lref{GlobalTheory.TwistAbove.twist_above}, which states that if
$r'(g)=\mu(g)r(g)$ for all $g$, then $r'$ is de Rham (crystalline, potentially crystalline) whenever $r$
and $\mu^{(1)}$ are, and that the jumps of $r'$ at $\tau$ are those of $r$ shifted by the jump of
$\mu^{(1)}$. (This is the compatibility of $D_{\mathrm{dR}}$ and $D_{\mathrm{cris}}$ with tensor
products, \cite{Fontaine1994}.)
\end{proof}

\begin{lemma}\label{pf:geo:wdfacts}
Let $A$ be a place of $\Omega$ over $w$ and $x,y$ Weil--Deligne representations of $W_A$ with $y$
one-dimensional. If $g\in W_A$, $\mathrm{cp}_g(x)=\geounr_S(\deg g)$ and $\mathrm{cp}_g(y)=\geounr_{\{c\}}(\deg g)$, then
$\mathrm{cp}_g(x\otimes y)=\geounr_{cS}(\deg g)$, where $cS=\{cz:z\in S\}$. Moreover
$\mathrm{cp}_g(x\oplus y')=\mathrm{cp}_g(x)\,\mathrm{cp}_g(y')$ for all $x,y'$; a direct summand of an unramified
Weil--Deligne representation is unramified; and the tensor product of two unramified ones is
unramified.
\formallabel{UpperTheory.Gal.TP.Sub.wDFacts_holds}
\end{lemma}

\begin{proof}
Since $y$ is one-dimensional, $y(g)$ is the scalar $c^{-\deg g}$, so $(x\otimes y)(g)=c^{-\deg g}x(g)$ and
the eigenvalues of $(x\otimes y)(g)$ are those of $x(g)$ multiplied by $c^{-\deg g}$; if they are the
$z^{-\deg g}$, $z\in S$, then they are the $(cz)^{-\deg g}$. The other assertions are immediate from
the block structure of direct sums and the definitions.
\end{proof}

\subsubsection{The twisting statements}

\begin{proposition}[(S7)--(S10)]\label{pf:geo:twists}
Let $\rho\colon G_E\to\GL_{dd}(\Qlbar)$ and $\chi\colon G_E\to\Qlbar^\times$ be continuous, $c\in\C$ and
$m\colon\Hom(E,\C)\to\Z$, and assume $\geoPar_w(\chi^{(1)};\{c\},\tau\mapsto\{m(\iota\circ\tau)\})$. Then:
\begin{enumerate}[label=\textup{(S\arabic*)},start=7,nosep]
\item if $\ell\notin w$, then $\geoES_w(\rho,R)$ implies $\geoES_w(\geotw(\rho,\chi),cR)$ for every multiset
  $R\subset\C$;
\item if $\rho$ is crystalline at $w$, so is $\geotw(\rho,\chi)$;
\item for every predicate $Q$ on $\Hom(E,\C)\times\Z$, $\geoJ_w(\rho,Q)$ implies $\geoJ_w(\geotw(\rho,\chi),Q')$
  with $Q'(\tau,j)=Q(\tau,j-m(\tau))$;
\item $\geoPhiK_w(\rho,R)$ implies $\geoPhiK_w(\geotw(\rho,\chi),cR)$ for every multiset $R\subset\C$.
\end{enumerate}
\formallabel{UpperTheory.Gal.Int.Geo.twES}\formallabel{UpperTheory.Gal.Int.Geo.twCrys}\formallabel{UpperTheory.Gal.Int.Geo.twJ}\formallabel{UpperTheory.Gal.Int.Geo.twPhi}
\end{proposition}

\begin{proof}
(S7) Let $A$ be a place over $w$. By $\geoPar_w$(a) the character $\chi$ is trivial on $I_A$, and so is
$\rho$ by $\geoES_w(\rho,R)$; hence $\geotw(\rho,\chi)$ is trivial on $I_A$. Let $\sigma$ be an arithmetic
Frobenius element. The characteristic polynomial of the $1\times1$ matrix $\chi(\sigma^{-1})$ is
$X-\chi(\sigma^{-1})$, so $\geoPar_w$(a) gives $\iota(\chi(\sigma^{-1}))=c$. Put $\mu=\iota^{-1}(c)$, $M=\rho(\sigma^{-1})$ and
$P=\prod_{r\in R}(X-\iota^{-1}r)$, so that $P(M)=0$. Then $\geotw(\rho,\chi)(\sigma^{-1})=\mu M$, and
\[
  \prod_{r\in R}\bigl(\mu M-\mu\,\iota^{-1}r\bigr)=\mu^{|R|}P(M)=0 .
\]
Since $\iota^{-1}(cR)=\mu\,\iota^{-1}(R)$, this is the required annihilation.

(S8) Fix a tied period package, a place $A$ over $w$ and a bridge inducing $A$. By $\geoPar_w$(b),
$\chi^{(1)}|_{D_A}$ is crystalline; $\rho|_{D_A}$ is crystalline by hypothesis; and
$\geotw(\rho,\chi)|_{D_A}=\geotw(\rho|_{D_A},\chi|_{D_A})$. Apply \cref{pf:geo:twloc}(a).

(S9) With the same data and an embedding $\tau$ induced by the bridge, $\geoJ_w(\rho,Q)$ gives a jump
multiset $J$ of $\rho|_{D_A}$ at $\tau$ containing every $j$ with $Q(\iota\circ\tau,j)$, and $\geoPar_w$(b) gives
the jump multiset $\{m(\iota\circ\tau)\}$ of $\chi^{(1)}|_{D_A}$. The restrictions are continuous, so by
\cref{pf:geo:twloc}(b) the multiset $J+m(\iota\circ\tau)$ is a jump multiset of $\geotw(\rho,\chi)|_{D_A}$ at $\tau$. If
$Q'(\iota\circ\tau,j)$, i.e.\ $Q(\iota\circ\tau,j-m(\iota\circ\tau))$, then $j-m(\iota\circ\tau)\in J$, so $j\in J+m(\iota\circ\tau)$.

(S10) With the same data, $\geoPar_w$(b) gives that $\chi^{(1)}|_{D_A}$ is crystalline with
$\mathrm{cp}_g(\geoWD(\chi^{(1)}|_{D_A}))=\geounr_{\{c\}}(\deg g)$ for all $g$, and $\geoPhiK_w(\rho,R)$ gives that
$\rho|_{D_A}$ is crystalline and that there are $g\in W_A$ with $\deg g=-1$ and $k\in\N$ with
$\mathrm{cp}_g(\geoWD(\rho|_{D_A}))\mid(\geort_R)^k$. By \cref{pf:geo:twloc}(a) the twist is crystalline. Crystalline
representations are potentially semistable, so by
\lref{GlobalTheory.T.fontaineParameter_twist_potentiallySemistable} (the Weil--Deligne
representation of the twist of a potentially semistable representation by a potentially semistable
character is the tensor product of their Weil--Deligne representations)
$\geoWD(\geotw(\rho,\chi)|_{D_A})=\geoWD(\rho|_{D_A})\otimes\geoWD(\chi^{(1)}|_{D_A})$. Write
$\mathrm{cp}_g(\geoWD(\rho|_{D_A}))=\geort_{S'}=\geounr_{S'}(-1)$ with $S'$ its multiset of roots. By
\cref{pf:geo:wdfacts}, $\mathrm{cp}_g(\geoWD(\geotw(\rho,\chi)|_{D_A}))=\geort_{cS'}$. Finally
$\geort_{S'}\mid(\geort_R)^k$ implies $\geort_{cS'}\mid(\geort_{cR})^k$: for $c\ne0$ substitute $X/c$ and multiply by
the appropriate powers of $c$; for $c=0$ both sides are powers of $X$ and the degrees compare as
before.
\end{proof}

\subsubsection{Normalisations}

\begin{lemma}\label{pf:geo:satpol}
Let $A$ be a place of $\Omega$ over $w$, $a\in\N$ and $B$ a multiset of nonzero complex numbers. For every
$g\in W_A$ the characteristic polynomial of $g$ on the Weil--Deligne representation which is unramified
with Frobenius eigenvalues $B^{-1}$ and twisted by $|\cdot|^{(1-a)/2}$ is $\geounr_{(\Nm w)^{(a-1)/2}B}(\deg g)$.
In particular for $\deg g=-1$ it is $\prod_{b\in B}(X-(\Nm w)^{(a-1)/2}b)$.
\formallabel{UpperTheory.Gal.TP.Sub.satPol_holds}
\end{lemma}

\begin{proof}
A direct computation from the definitions of unramified Weil--Deligne representations and of the twist
by $|\cdot|^s$: an element of degree $n$ acts on the line of $b\in B$ by $(\Nm w^{(a-1)/2}b)^{-n}$.
\end{proof}

\begin{proposition}[(S11)--(S13)]\label{pf:geo:norms}
\begin{enumerate}[label=\textup{(S\arabic*)},start=11,nosep]
\item If $B$ is a multiset of $a$ nonzero complex numbers and $\geoPhiK_w(\rho,(\Nm w)^{(a-1)/2}B)$ holds, then
  $\geoPhiK^{\mathrm{sat}}_w(\rho;a,B)$ holds.
\item If $E$ is a CM field and $B$ is a Satake multiset of $U$ at $w$, then $B$ has $a$ elements.
\item Let $E$ be a CM field, $U$ a space of functions on $\geoUa$, $\rho\colon G_E\to\GL_{dd}(\Qlbar)$ and $Q$ a
  predicate on $\Hom(E,\C)\times\Z$ such that $Q(\tau,j)$ holds for every $\tau$ over $w$, every weight $\lambda$
  of $U$ at $\tau$ and every $j\in J(\lambda)$. Then $\geoJ_w(\rho,Q)$ implies $\geoW_w(U,\rho)$.
\end{enumerate}
\formallabel{UpperTheory.Gal.Int.Geo.phiOf}\formallabel{UpperTheory.Gal.Int.Geo.satCard}\formallabel{UpperTheory.Gal.Int.Geo.jumpsOf}
\end{proposition}

\begin{proof}
(S11) The two conditions differ only in the polynomial: for the element $g$ of degree $-1$ provided by
$\geoPhiK_w$, the polynomial $\geort_{(\Nm w)^{(a-1)/2}B}$ is the characteristic polynomial of $g$ on the
twisted unramified Weil--Deligne representation, by \cref{pf:geo:satpol}.

(S12) By definition $B$ is the Satake multiset of a spherical smooth irreducible representation of
$\GL_a(E_w)$; such multisets have $a$ elements (\lref{GlobalTheory.LFun.HasSatGL.card}: the Satake
parameters of a spherical representation of $\GL_a$ form a multiset of cardinality $a$).

(S13) Fix the auxiliary data of the definition of $\geoW_w(U,\rho)$: a tied period package, a place
$A$ over $w$, a bridge inducing $A$, an embedding $\tau_0\colon E\to\Qlbar$ induced by it, a weight $\lambda$ of
$U$ at $\iota\circ\tau_0$ and a multiset $J$ of integers whose image in $\C$ is
$\{\frac{a-1}2-z:z\in\inf(\lambda)\}$, $\inf(\lambda)=\{\lambda_k+\frac{a-1}2-k\}_k$. By \cref{pf:geo:wtjumps} below,
$J=J(\lambda)$. By \cref{pf:geo:overpl}, $\iota\circ\tau_0$ lies over $w$, so $Q(\iota\circ\tau_0,j)$ holds for all $j\in J$,
and $\geoJ_w(\rho,Q)$ provides a jump multiset of $\rho|_{D_A}$ at $\tau_0$ containing $J$.
\end{proof}

\begin{lemma}\label{pf:geo:wtjumps}
Let $\lambda=(\lambda_0\ge\dots\ge\lambda_{a-1})\in\Z^a$. The map $k\mapsto k-\lambda_k$ is strictly increasing, so
$J(\lambda)$ has $a$ elements, and the image of the multiset $J(\lambda)$ in $\C$ is
$\{\frac{a-1}2-z:z\in\inf(\lambda)\}$. Consequently a multiset of integers with this image equals
$J(\lambda)$.
\formallabel{GlobalTheory.RoadWGeo.strictMono_jump}\formallabel{GlobalTheory.RoadWGeo.card_wtJumps}\formallabel{GlobalTheory.RoadWGeo.wtJumps_val}
\end{lemma}

\begin{proof}
If $i<j$ then $\lambda_i\ge\lambda_j$, so $i-\lambda_i<j-\lambda_j$. For each $k$,
$\frac{a-1}2-(\lambda_k+\frac{a-1}2-k)=k-\lambda_k$. The last assertion holds because the map $\Z\to\C$ is injective.
\end{proof}

\begin{lemma}\label{pf:geo:overpl}
Let $A$ be a place of $\Omega$ over $w$, $B_w$ a bridge inducing $A$, and $\tau_0\colon E\to\Qlbar$ induced by
$B_w$. Then $\iota\circ\tau_0$ lies over $w$, i.e.\ $|\tau_0(y)|<1$ for every $y\in w$.
\formallabel{GlobalTheory.CongL.LimAtLU.overPl_of_induced}
\end{lemma}

\begin{proof}
Write $j\colon\Omega\to\Qlbar$ for the embedding of the bridge, so that $A=\{z:|j(z)|\le1\}$, and
$\tau_0=\tau'\circ j$ on $E$ for a $\Ql$-embedding $\tau'$ of the local field $K_w$ of the bridge. For $y\in w$
we have $y\in\mathfrak m_A$, hence $|j(y)|<1$, and $\Ql$-embeddings of finite extensions preserve the
absolute value, so $|\tau_0(y)|=|j(y)|<1$.
\end{proof}

\subsubsection{The $\ell$-adic characters of algebraic Hecke characters}

A Hecke character $\chi$ of $E$ is \emph{algebraic of infinity type $n\colon\Hom(E,\C)\to\Z$} if
$\chi_\infty(k)=\prod_\phi\phi(k)^{n(\phi)}$ for every $k\in E^\times$ which is positive at all real places. A global
Artin map is fixed as in \cref{def:galois:wd}. A character $r\colon G_E\to\Qlbar^\times$ is \emph{an $\ell$-adic
character of $(\chi,n)$} if
\begin{equation}\label{eq:pf:geo:ladic}
  r(\sigma)=\iota^{-1}\bigl(\chi(x)\,\chi_n(x_\infty)^{-1}\bigr)\cdot\prod_{\phi}\phi_\ell(x_\ell)^{n(\phi)}
\end{equation}
whenever $\mathrm{art}(x)$ is the image of $\sigma$ in $G_E^{\mathrm{ab}}$; here
$\chi_n(x_\infty)=\prod_\phi\phi(x_\infty)^{n(\phi)}$ and $\phi_\ell$ is the $\ell$-adic embedding attached to $\phi$
through $\iota$. This is Weil's and Serre's construction \cite{Weil1956,Serre1968}.

\begin{littheorem}[{\cite{Weil1956}; \cite[Ch.~II, III]{Serre1968}; \cite{Fontaine1994}}]\label{pf:geo:lit:ladic}
\begin{enumerate}[label=(\alph*),nosep]
\item For every number field $E$ there is a global Artin map such that for every prime $\ell$, every
  $\iota$ and every Hecke character $\chi$ algebraic of infinity type $n$ there is a unique continuous
  $\ell$-adic character of $(\chi,n)$.
\item Such a character is de Rham at every place above $\ell$ (it is locally algebraic), and for every
  embedding $\tau$ induced by a bridge its jump multiset at $\tau$ is $\{-n(\iota\circ\tau)\}$ (with the conventions
  for jumps of the development).
\end{enumerate}
\formallabel{PositivityTheory.Stage1.PM.heckeLAdicCharacter}\formallabel{HodgeLocalDeRhamField.HLdR.field_outright}
\end{littheorem}

Both statements are labelled as published results in the development, where they are proved (the
first from global Artin reciprocity). The theorem that de Rham representations are potentially
semistable \cite{Berger2002} is likewise a literature statement proved in the development.

\begin{proposition}[(S15)]\label{pf:geo:charex}
For every number field $E$, $\Omega$, $\ell$ and $\iota$ there is a map $\mathcal L$ from pairs $(\chi,n)$ to
characters $G_E\to\Qlbar^\times$ such that, for $\chi$ algebraic of infinity type $n$, $\mathcal L(\chi,n)$ is
continuous and $\geoPar_w(\mathcal L(\chi,n)^{(1)};\{\chi(w)\},\tau\mapsto\{-n(\iota\circ\tau)\})$ holds at every place
$w$ at which $\chi$ is unramified, $\chi(w)$ denoting the value of $\chi$ at a uniformizer of $E_w$; and
$\mathcal L(\chi,n)=\mathcal L(\chi',n')$ implies $\chi=\chi'$ for algebraic $(\chi,n)$, $(\chi',n')$.
\formallabel{UpperTheory.Gal.TP.Sub.charEx}\formallabel{UpperTheory.Gal.TP.Sub.charFrob_holds}\formallabel{UpperTheory.Gal.TP.Sub.charAbove_holds}\formallabel{UpperTheory.Gal.TP.Sub.charInj_holds}
\end{proposition}

\begin{proof}
Let $\mathrm{art}$ be the global Artin map of \cref{pf:geo:lit:ladic} and let $\mathcal L(\chi,n)$ be the unique
continuous $\ell$-adic character of $(\chi,n)$ if $\chi$ is algebraic of infinity type $n$, and $1$ otherwise.
Let $\chi$ be algebraic of type $n$ and unramified at $w$, and put $r=\mathcal L(\chi,n)$.

\emph{Away from $\ell$.} Let $\ell\notin w$ and $A$ over $w$. For $x$ in the local units at $w$
(embedded in the ideles), \eqref{eq:pf:geo:ladic} gives $r(\sigma)=\iota^{-1}\chi_w(x)=1$ for $\sigma$ with
$\mathrm{art}(x)=\sigma^{\mathrm{ab}}$, since the factors at $\infty$ and $\ell$ are trivial; these $\sigma$ are the
elements of the inertia groups over $w$, so $r$ is trivial on $I_A$. A uniformizer $\varpi_w$ is mapped by
$\mathrm{art}$ to a geometric Frobenius, so for an arithmetic Frobenius $\sigma$ one gets
$r(\sigma)=\iota^{-1}(\chi(w))^{-1}$, i.e.\ the characteristic polynomial of $r(\sigma^{-1})$ is $X-\iota^{-1}\chi(w)$.
In the development these two facts are \lref{GlobalTheory.HeckeGal.IsLAdicOf.unr_iff} ($r$ is trivial
on the inertia groups over $w\nmid\ell$ if and only if $\chi$ is unramified at $w$) and
\lref{GlobalTheory.HeckeGal.IsLAdicOf.frob_unr} (the value of $r$ at arithmetic Frobenius elements over
such $w$ is $\iota^{-1}(\chi(w))^{-1}$). This is $\geoPar_w$(a).

\emph{Above $\ell$.} Fix a tied period package, a place $A$ over $w$ and a bridge inducing $A$. By
\cref{pf:geo:lit:ladic}(b), $r^{(1)}|_{D_A}$ is de Rham with jump multiset $\{-n(\iota\circ\tau)\}$ at every
induced $\tau$; it is continuous, hence potentially semistable \cite{Berger2002}. By
\lref{UpperTheory.Lit.CharWD.algebraicCharacterWD} its Weil--Deligne representation is the character
$\chi_w\circ a$ of $W_A$, where $a$ is the local Artin map at $A$ (\cref{def:galois:wd}) and $\chi_w$ the local
component of $\chi$. Since $\chi$ is unramified at $w$, this Weil--Deligne representation is unramified, so
$r^{(1)}|_{D_A}$ is crystalline (\lref{GlobalTheory.PstLaws.unr_wdAt_iff}: a potentially semistable
representation is crystalline if and only if its Weil--Deligne representation is unramified); and
$\chi_w(a(g))=\chi(w)^{-\deg g}$ because $\chi_w(u)=\chi(w)^{\ord_w(u)}$ and $\ord_w(a(g))=-\deg g$. Hence
$\mathrm{cp}_g=X-\chi(w)^{-\deg g}=\geounr_{\{\chi(w)\}}(\deg g)$. This is $\geoPar_w$(b).

\emph{Injectivity.} If $r$ is an $\ell$-adic character of
both $(\chi,n)$ and $(\chi',n')$, then at every place $v\nmid\ell$ where both $\chi,\chi'$ are unramified,
applying \eqref{eq:pf:geo:ladic} to the idele $\varpi_v$ gives $\iota^{-1}\chi(v)=r(\sigma)^{-1}=\iota^{-1}\chi'(v)$.
Two Hecke characters which are unramified and agree at uniformizers at all but finitely many places
are equal (\lref{GlobalTheory.T.hecke_character_eq_of_cofinite}).
\end{proof}

\subsubsection{Semisimplicity and continuity}

\begin{proposition}[(S16), (S17)]\label{pf:geo:ssc}
Let $\Gamma$ be a group, $k$ a field, $r\colon\Gamma\to\GL_a(k)$, $s\colon\Gamma\to\GL_b(k)$ and
$c\colon\Gamma\to k^\times$ homomorphisms, $r\oplus s$ the block diagonal sum, and $r'$ a conjugate of $r$
($r'=uru^{-1}$ with $u\in\GL_a(k)$).
\begin{enumerate}[label=\textup{(S\arabic*)},start=16,nosep]
\item If $r\oplus s$ is semisimple, so is $r$; if $r$ is semisimple, so are $r'$ and $\geotw(r,c)$.
\item Assume $\Gamma$ is a topological group and $k$ a topological field. If $r\oplus s$ is continuous, so is
  $r$; if $r$ is continuous, so is $r'$; if $r$ and $c$ are continuous, so are $\geotw(r,c)$ and $c^{(1)}$.
\end{enumerate}
Moreover the characteristic polynomial of $(r\oplus s)(g)$ is the product of those of $r(g)$ and $s(g)$,
that of $c^{(1)}(g)$ is $X-c(g)$, and conjugate representations have the same characteristic polynomials.
\formallabel{UpperTheory.Gal.TP.Sub.sSLeft_holds}\formallabel{UpperTheory.Gal.TP.Sub.contLeft_holds}\formallabel{UpperTheory.Gal.TP.Sub.cpFacts_holds}
\end{proposition}

\begin{proof}
(S16) Let $W\subseteq k^a$ be $r$-stable. Then $W'=W\oplus k^b$ is stable under $r\oplus s$; let $C$ be a
stable complement. The projection $\pi\colon k^{a+b}\to k^a$ is equivariant, so $\pi(C)$ is $r$-stable, and
$\pi(C)+W=\pi(C+W')=k^a$. If $\pi(c)=w\in W$ with $c\in C$, then $c-(w,0)\in0\oplus k^b\subseteq W'$, so $c\in W'\cap C=0$;
hence $\pi(C)\cap W=0$ and $\pi(C)$ is a stable complement of $W$. The stable subspaces of $r'$ are the
images under $u$ of those of $r$, and $\geotw(r,c)$ has the same stable subspaces as $r$ since $c(g)$ is a
scalar. (S17) The entries of $r(g)$ are entries of $(r\oplus s)(g)$; $r'(g)=ur(g)u^{-1}$ and
$\geotw(r,c)(g)=c(g)r(g)$ are continuous in $g$ since multiplication is continuous; $c^{(1)}$ is $c$. The
assertions on characteristic polynomials are standard.
\end{proof}

% ===========================================================================
\subsection{Facts from $p$-adic Hodge theory used by the steps}\label{sec:pf:geo:padic}

The following classical facts are used in the proofs of (S3)--(S5) in \cref{sec:pf:steps1}. They are
proved in the formal development; we state the forms used and cite the literature. Notation for
period packages, bridges and jump multisets is that of \cref{def:pf:steps2:crys}; for a finite
extension $K/\Ql$, $G_K$ is its absolute Galois group, $\chi_{\mathrm{cyc}}$ the cyclotomic character and
$\C_\ell$ the completion of $\Qlbar$.

\begin{proposition}[Hodge--Tate decomposition]\label{pf:geo:htbasis}
Let $K,LL$ be finite extensions of $\Ql$ such that every $\Ql$-embedding of $K$ takes values in $LL$, let
$\tau'\colon K\to LL$ be the inclusion, $m\in\N$, and $r'\colon G_K\to\GL_m(LL)$ a homomorphism which is de Rham
in the sense that $\dim_{LL}D_{\mathrm{dR}}(r')=[K:\Ql]\,m$ (for a tied period package). Let $K'$ be a finite
extension containing $K$ and $LL$. Then there are $P\in\GL_m(\C_\ell)$ and integers $k_1,\dots,k_m$ such that
each $j\in\Z$ occurs among the $k_i$ with multiplicity equal to the $\tau'$-Hodge--Tate multiplicity of $j$ (the
dimension of the $j$-th graded piece of the Hodge filtration on the $\tau'$-component of $D_{\mathrm{dR}}$),
and
\[
  r'(g)\,{}^gP=P\,\mathrm{diag}\bigl(\chi_{\mathrm{cyc}}(g)^{-k_1},\dots,\chi_{\mathrm{cyc}}(g)^{-k_m}\bigr)\qquad(g\in G_{K'}),
\]
where ${}^gP$ is obtained by applying $g$ to the entries of $P$.
\formallabel{GlobalTheory.T.deRham_exists_hodgeTate_eigenbasis}
\end{proposition}

\begin{proof}
De Rham representations are Hodge--Tate, and a Hodge--Tate representation becomes, after extension of
scalars to $\C_\ell$, a direct sum of the twists $\C_\ell(-j)$ with multiplicities given by the Hodge--Tate
numbers \cite{Fontaine1994}, \cite{Tate1967pdiv}; the columns of $P$ are a basis adapted to this
decomposition for the $\tau'$-component. The formal development proves this classical statement.
\end{proof}

\begin{proposition}\label{pf:geo:pstfacts}
Fix a tied period package.
\begin{enumerate}[label=(\alph*),nosep]
\item A crystalline representation of $G_K$ on a finite-dimensional $LL$-vector space is semistable, and a
  semistable one is potentially semistable.
\item The Hodge--Tate multiplicities $\dim\mathrm{gr}^jD_{\mathrm{dR}}$ at an embedding are unchanged when
  the bridge is replaced by a conjugate bridge and the embedding by the corresponding conjugate embedding.
\item If $r$ is a representation of a decomposition group $D_A$ with a model $r'$ over a finite extension
  $LL$ which is potentially semistable, then the characteristic polynomial of every $g\in W_A$ on the
  Weil--Deligne representation of $r$ is that of $g$ on Fontaine's Weil--Deligne representation of $r'$
  \cite{Fontaine1994rss}, and this Weil--Deligne representation has dimension equal to the degree of $r$.
\end{enumerate}
\formallabel{GlobalTheory.CrysSst.isSst_of_isCrys}\formallabel{GlobalTheory.GalCut.IwaTwist.isPst_of_isSst}\formallabel{GlobalTheory.Pst.BridgeInd.multV_conj}\formallabel{UpperTheory.Gal.crystallineComparison.wdModel_of}
\end{proposition}

\begin{proof}
(a) $D_{\mathrm{cris}}(V)$ embeds into $D_{\mathrm{st}}(V)$, and $\dim_{K_0}D_{\mathrm{st}}(V)\le\dim V$; if $V$ is
crystalline the first space has dimension $\dim V$, hence so does the second \cite{Fontaine1994rss}. A
semistable representation is potentially semistable (take the trivial extension). (b) Conjugation by an
element of $G_{\Ql}$ transports the Hodge filtration of one bridge to that of the conjugate one,
preserving dimensions. (c) Fontaine's Weil--Deligne representation does not depend on the choice of the
finite extension over which a potentially semistable model is chosen (enlarging the coefficients does not
change it), and it has the dimension of the representation.
\end{proof}

\begin{lemma}\label{pf:geo:fetbasic}
Let $d,e\in\N$. A base change of a morphism which is finite \'etale of degree $d$ is finite \'etale of degree
$d$; the composite of morphisms which are finite \'etale of degrees $e$ and $d$ is finite \'etale of degree
$de$; and a Galois cover with finite group $G$ (\cref{def:pf:steps2:hop}) is finite \'etale of degree $|G|$.
\formallabel{UpperTheory.Gal.G3.fEtBasic_holds}
\end{lemma}

\begin{proof}
Finiteness and \'etaleness are stable under base change and composition. Geometric points of a base change
over a geometric point correspond to geometric points of the original fibre, and in a composite the lifts of a
geometric point are the lifts, through the second map, of each of its lifts through the first. For a Galois
cover $W\to X$ with group $G$, the isomorphism $\coprod_GW\cong W\times_XW$ shows that $G$ acts simply
transitively on the geometric points of each nonempty fibre; fibres are nonempty since Galois covers are
surjective.
\end{proof}

% ===========================================================================
\subsection{Hodge cohomology, idempotents, Schur polynomials and class numbers}\label{sec:pf:geo:hodge}

This subsection collects results used in the construction of the Kuga--Sato package
(\cref{sec:pf:ks1,sec:pf:ks2}).

\subsubsection{Hodge cohomology}

\begin{definition}[Hodge cohomology, pull-back and trace]\label{pf:geo:hodgedef}
Let $k$ be a field and $s\colon X\to\Spec k$ a scheme over $k$. For $j\in\N$, $\Omega^j_{X/k}$ is the sheaf of
$\OO_X$-modules associated with the presheaf $V\mapsto\bigwedge^j\Omega_{\OO(V)/k}$, and
$\geoHdg{X}{j}{b}=H^b(X,\Omega^j_{X/k})$ is its sheaf cohomology (the derived-functor cohomology of the
underlying sheaf of abelian groups, a $k$-vector space). For a morphism $f\colon X'\to X$ over $k$, the
pull-back $f^*\colon\geoHdg{X}{j}{b}\to\geoHdg{X'}{j}{b}$ is induced by the pull-back of differential forms
$\Omega^j_{X/k}\to f_*\Omega^j_{X'/k}$. If $f$ is finite and \'etale, then on every affine open $V\subseteq X$
the ring $\OO(f^{-1}V)$ is finite \'etale over $\OO(V)$, so $\Omega^j_{\OO(f^{-1}V)}=\OO(f^{-1}V)\otimes_{\OO(V)}
\Omega^j_{\OO(V)}$, and the maps $\mathrm{Tr}_{\OO(f^{-1}V)/\OO(V)}\otimes\mathrm{id}$ glue to the \emph{trace of forms}
$f_*\Omega^j_{X'/k}\to\Omega^j_{X/k}$; the trace $f_*\colon\geoHdg{X'}{j}{b}\to\geoHdg{X}{j}{b}$ is the map it
induces on cohomology, composed with the identification $H^b(X',\Omega^j_{X'/k})\cong H^b(X,f_*\Omega^j_{X'/k})$
for the affine morphism $f$.
\formal{GlobalTheory.HodgeCoh.Omega}\formal{GlobalTheory.HodgeCoh.hodgeH}\formal{GlobalTheory.HodgeCoh.hodgePull}\formal{GlobalTheory.HodgeCoh.hodgeTrace}
\end{definition}

Pull-backs are functorial, and $f_*\circ f^*$ is induced by the composite $\Omega^j_X\to f_*\Omega^j_{X'}\to\Omega^j_X$ of
the pull-back and the trace of forms. For a morphism $s\colon Y\to Z$, a morphism $c\colon W\to Z$ and a finite
family of morphisms $\sigma_i\colon W\to Y$ with $s\sigma_i=c$, we say that the $\sigma_i$ \emph{split $s$ over $c$} if
the induced morphism $\coprod_iW\to Y\times_ZW$ is an isomorphism (\cref{def:pf:steps2:hop}).

\begin{lemma}\label{pf:geo:hodgeinj}
Let $k$ be a field of characteristic $0$ and $c\colon W\to Z$ a finite, \'etale and surjective morphism of
$k$-schemes. Then $c^*\colon\geoHdg{Z}{a}{b}\to\geoHdg{W}{a}{b}$ is injective for all $a,b$.
\formallabel{UpperTheory.Gal.G3.pHodgeInj}
\end{lemma}

\begin{proof}
The composite $\Omega^a_Z\to c_*\Omega^a_W\to\Omega^a_Z$ of pull-back and trace of forms is multiplication by the
degree of $c$, a locally constant function which is a positive integer everywhere because $c$ is
surjective; since $k$ has characteristic $0$ it is invertible, so there is an endomorphism $\psi$ of $\Omega^a_Z$
with $\psi\circ(\text{trace}\circ\text{pull-back})=\mathrm{id}$. On cohomology, $c_*c^*$ is induced by the composite,
so $\psi_*\circ c_*\circ c^*=\mathrm{id}$ and $c^*$ is injective.
\end{proof}

\begin{lemma}\label{pf:geo:hodgetr}
Let $k$ be a field and $s\colon Y\to Z$ a finite \'etale morphism of $k$-schemes. Assume that the trace of
functions is compatible with splittings: for every $c\colon W\to Z$, every family $(\sigma_i)$ splitting $s$ over $c$,
every affine open $V\subseteq Z$ and every $b\in\OO(s^{-1}V)$ one has
$c^*(\mathrm{Tr}(b))=\sum_i\sigma_i^*(b)$ in $\OO(c^{-1}V)$. Then the trace on Hodge cohomology satisfies
$c^*(s_*x)=\sum_i\sigma_i^*x$ for all $c$, all splitting families $(\sigma_i)$ and all $x\in\geoHdg{Y}{a}{b}$.
Consequently, if $k$ has characteristic $0$, the theory $X\mapsto\geoHdg{X}{a}{b}$ is \emph{good}: pull-back
along Galois covers is injective, and every finite \'etale morphism has a trace in the above sense.
\formallabel{UpperTheory.Gal.G3.pHodgeTr}\formallabel{UpperTheory.Gal.G3.pHodgeGood}\formallabel{GlobalTheory.GalTrace.omegaTrace_omegaPull}\formallabel{GlobalTheory.GalTrace.hodgePull_hodgeTrace}
\end{lemma}

\begin{proof}
On an affine open $V$, a form on $s^{-1}V$ is a sum of terms $b\,s^*\eta$ with $b\in\OO(s^{-1}V)$ and $\eta$ a form on
$V$, and the trace of forms sends it to $\mathrm{Tr}(b)\eta$; the hypothesis gives
$c^*(\mathrm{Tr}(b)\eta)=\sum_i\sigma_i^*(b)\,c^*\eta=\sum_i\sigma_i^*(b\,s^*\eta)$, because $\sigma_i^*s^*=c^*$. Hence the
morphisms of sheaves $c^*\circ\mathrm{trace}$ and $\sum_i\sigma_i^*$ from $s_*\Omega^a_Y$ to $c_*\Omega^a_W$ agree on all
affine opens, hence agree, and the identity passes to cohomology by functoriality of the maps induced on
sheaf cohomology. For the last assertion, Galois covers are finite, \'etale and surjective, so
\cref{pf:geo:hodgeinj} applies; and the hypothesis on traces of functions holds for every finite \'etale $s$
(\lref{UpperTheory.Gal.G3.hodgeTr_holds}, which applies the first part to a verification of this
hypothesis).
\end{proof}

\begin{lemma}[Operators of correspondences]\label{pf:geo:optr}
Let $T$ be a good cohomology theory as in \cref{pf:geo:hodgetr} (for instance Hodge cohomology over a field of
characteristic $0$), $Z\xleftarrow{s}Y\xrightarrow{t}Z'$ morphisms with $s$ finite \'etale, and $\mathrm{tr}$ a trace for $s$.
If $A\colon T(Z')\to T(Z)$ is an operator of the correspondence $(s,t)$ (there are a Galois cover $c\colon W\to Z$
and a family $(\sigma_i)$ splitting $s$ over $c$ with $c^*(Ax)=\sum_i(t\sigma_i)^*x$ for all $x$), then
$A=\mathrm{tr}\circ t^*$. Conversely, if $s$ splits over some Galois cover, then $\mathrm{tr}\circ t^*$ is an operator
of $(s,t)$.
\formallabel{UpperTheory.Gal.G3.pOpTr}
\end{lemma}

\begin{proof}
For a Galois cover $c$ and a splitting family $(\sigma_i)$,
$c^*(\mathrm{tr}(t^*x))=\sum_i\sigma_i^*t^*x=\sum_i(t\sigma_i)^*x$ by the trace property and functoriality. If $A$ is an
operator of $(s,t)$ with respect to $c,(\sigma_i)$, then $c^*(Ax)=c^*(\mathrm{tr}(t^*x))$, and $c^*$ is injective. The
converse is the same identity.
\end{proof}

\subsubsection{Idempotent certificates}

For a family $(T_j)_{j\in J}$ of elements of an $F$-algebra $A$, a word $u=(j_1,\dots,j_r)$ in $J$ and
$\theta\colon J\to F$, put $T_u=T_{j_1}\cdots T_{j_r}$ and $\theta_u=\theta_{j_1}\cdots\theta_{j_r}$; for a finitely supported
function $t$ on words put $t(T)=\sum_ut(u)T_u$ and $t(\theta)=\sum_ut(u)\theta_u$.

\begin{lemma}\label{pf:geo:projone}
Let $F$ be a field, $A$ a finite-dimensional $F$-algebra, $(T_j)_{j\in J}$ pairwise commuting elements of $A$ and
$\theta\colon J\to F$. There is a finitely supported $t$ such that $t(\theta)=1$, $t(T)$ is idempotent, and for every
$j$ there is $N$ with $(T_j-\theta_j)^N\,t(T)=0$.
\formallabel{UpperTheory.Gal.CertificateRingHomMap.pProjOne_holds}
\end{lemma}

\begin{proof}
Let $D\subseteq F\times A$ be the $F$-subalgebra generated by the pairs $(\theta_j,T_j)$; it is commutative and
finite-dimensional, and $\chi\colon D\to F$, $(x,y)\mapsto x$, is an $F$-algebra map. Let $b_1,\dots,b_r$ be a basis of $D$.
For each $i$ choose a polynomial $q_i$ such that $e_i=q_i(b_i)$ is idempotent, $\chi(e_i)=1$ and
$(b_i-\chi(b_i))e_i$ is nilpotent (the idempotent of the generalized eigenspace of $b_i$ for the eigenvalue
$\chi(b_i)$, which is a polynomial in $b_i$ since $b_i$ is integral over $F$). Then $z=\prod_ie_i$ is idempotent,
$\chi(z)=1$, and $(d-\chi(d))z$ is nilpotent for every $d\in D$: the set of such $d$ is an $F$-subspace containing
the $b_i$. Since $D$ is generated by the pairs $(\theta_j,T_j)$, $z=\sum_ut(u)(\theta_u,T_u)$ for some finitely supported
$t$; its components give $t(\theta)=\chi(z)=1$, $t(T)$ idempotent, and $(T_j-\theta_j)t(T)$ nilpotent, hence
$(T_j-\theta_j)^Nt(T)=((T_j-\theta_j)t(T))^N=0$ for large $N$ (the elements commute and $t(T)$ is idempotent).
\end{proof}

\subsubsection{Schur polynomials}

For $n\in\N$, $\lambda\in\N^n$ and elements $g_1,\dots,g_n$ of a commutative $\Q$-algebra, let
$s_\lambda(g)=\det\bigl(h_{\lambda_i-i+j}(g)\bigr)_{1\le i,j\le n}$ (Jacobi--Trudi; $h_k$ is the complete homogeneous
symmetric polynomial of degree $k$, $h_0=1$, $h_k=0$ for $k<0$), $V(g)=\prod_{i<j}(g_i-g_j)$ and
$a_e(g)=\det(g_j^{e_i})_{i,j}$ for $e\in\N^n$.

\begin{lemma}\label{pf:geo:schur}
Let $R$ be a commutative $\Q$-algebra, $n\in\N$, $\lambda\in\N^n$ and $g\in R^n$.
\begin{enumerate}[label=(\alph*),nosep]
\item (Bialternant formula) $s_\lambda(g)\,V(g)=a_{\lambda+\delta}(g)$, where $\delta=(n-1,n-2,\dots,0)$.
\item $s_0(g)=1$ and $s_{\lambda+(m,\dots,m)}(g)=(g_1\cdots g_n)^m s_\lambda(g)$ for $m\in\N$.
\item $s_\lambda(yg)=y^{|\lambda|}s_\lambda(g)$ for $y\in R$, $|\lambda|=\sum_i\lambda_i$.
\item $s_\lambda$ is symmetric; hence for a number field $F$, $\lambda\colon\Hom(F,\C)\to\N^n$ and $\chi\in F[X]$ whose images
  $\tau(\chi)$ have $n$ roots $g_\tau$ (with multiplicity), the product $\prod_\tau s_{\lambda(\tau)}(g_\tau)$ does not depend on
  the enumeration of the roots.
\end{enumerate}
\formallabel{UpperTheory.Gal.Schur.tie}\formallabel{UpperTheory.Gal.Schur.SchurDet_holds}\formallabel{UpperTheory.Gal.Schur.schurHomog}\formallabel{UpperTheory.Gal.LocalSmoothExtension.schWEq_of}
\end{lemma}

\begin{proof}
(a) is the classical equality of the Jacobi--Trudi determinant and the quotient of alternants
\cite[Ch.~I, (3.1), (3.4)]{Macdonald1995}; it is an identity of polynomials with rational coefficients, which
the development proves in the polynomial ring and specializes. (b) The matrix $(h_{-i+j})$ is upper
unitriangular, so $s_0=1$. In the polynomial ring $\Q[X_1,\dots,X_n]$, factoring $X_j^m$ out of the $j$-th column
gives $a_{\lambda+m+\delta}=(X_1\cdots X_n)^ma_{\lambda+\delta}$; by (a) and since $V(X)$ is a nonzerodivisor,
$s_{\lambda+m}=(X_1\cdots X_n)^ms_\lambda$; specialize to $g$. (c) $s_\lambda$ is homogeneous of degree $|\lambda|$, each
term of the determinant having degree $\sum_i(\lambda_i-i+\pi(i))=|\lambda|$. (d) $s_\lambda$ is a polynomial in the
symmetric functions $h_k$.
\end{proof}

\begin{lemma}\label{pf:geo:invroots}
Let $F$ be a number field, $P\in F[X]$ with $P(0)\ne0$, $P^\flat=P(0)^{-1}X^{\deg P}P(1/X)$ and $\tau\colon F\to\C$. The
multiset of roots of $\tau(P^\flat)$ is the multiset of inverses of the roots of $\tau(P)$.
\formallabel{UpperTheory.Gal.LocalSmoothExtension.InvRoots_holds}
\end{lemma}

\begin{proof}
If $\tau(P)=c\prod_i(X-r_i)$ with all $r_i\ne0$, then $X^{\deg P}\tau(P)(1/X)=c\prod_i(1-r_iX)=c\prod_i(-r_i)\prod_i(X-r_i^{-1})$
and $\tau(P(0))=c\prod_i(-r_i)$.
\end{proof}

\subsubsection{Level groups and class numbers}

Let $p$ be a prime and $D,G,L$ a datum, a unitary similitude group and a link for $p$
(\cref{pf:ks2:def:datum}). Recall that $K_N$ is the group of $g\in G(\A_f^p)$ such that $g$ and $g^{-1}$ act
trivially on $\hat\Lambda^p/N\hat\Lambda^p$, $\hat\Lambda^p=\hat\Z^p\otimes\Lambda$.

\begin{lemma}\label{pf:geo:kn}
For every $N\ge1$, $K_N$ is a compact open subgroup of $G(\A_f^p)$.
\formallabel{GlobalTheory.LKGeo.KNLevel.exists_level_eq_KN}
\end{lemma}

\begin{proof}
Choose a $\Z$-basis of $\Lambda$ and let $M(g)$ be the matrix of $g$ acting on $\A_f^p\otimes\Lambda$. Then $g\in K_N$ if and only
if $N^{-1}(M(g)-1)$ and $N^{-1}(M(g^{-1})-1)$ have entries in $\hat\Z^p$. The maps $g\mapsto N^{-1}(M(g^{\pm1})-1)$ are
continuous and the set of integral matrices is open, so $K_N$ contains an open neighbourhood of $1$ and is an
open, hence closed, subgroup. An element of $G(\A_f^p)$ is a pair $(x,\nu)$ with $xx^*=\nu$; if $x$ and $x^{-1}$ are
integral, then $x,x^{-1}$ lie in the compact set $P$ of integral elements of $\A_f^p\otimes C$ and $\nu,\nu^{-1}$ lie in the
compact image of $P$ under $y\mapsto yy^*$ (read in the centre). Since $G(\A_f^p)$ is closed in the product of the
unit groups, embedded as usual by $u\mapsto(u,u^{-1})$, the elements $g$ with $g,g^{-1}$ integral form a compact set;
it contains $K_N$, which is closed, so $K_N$ is compact.
\end{proof}

A \emph{group locus} over a number field $F$ is a set $S$ of polynomials in the matrix entries such that for every
commutative $F$-algebra $R$ the set $\mathcal G(R)$ of $g\in\GL_n(R)$ annihilating $S$ is a subgroup; i.e.\ $\mathcal G$ is a
closed subgroup scheme of $\GL_{n,F}$ given by explicit equations.

\begin{theorem}[Finiteness of class numbers]\label{pf:geo:classnumber}
Let $F$ be a number field and $\mathcal G\subseteq\GL_{n,F}$ a group locus. For every compact open subgroup $U$ of
$\mathcal G(\A_{F,f})$ there is a finite set $T\subseteq\mathcal G(\A_{F,f})$ such that every $x\in\mathcal G(\A_{F,f})$ can be written
$x=\gamma su$ with $\gamma\in\mathcal G(F)$, $s\in T$ and $u\in U$.
\formallabel{GlobalTheory.classNumberFinite_holds}\formallabel{GlobalTheory.T.isCompact_integralPoints_finiteAdelic_algebraicGroup}\formallabel{GlobalTheory.T.compact_meets_finitely_many_leftCosets_open}
\end{theorem}

\begin{proof}
For the integral level $U_0=\mathcal G(\A_{F,f})\cap\GL_n(\hat\OO_F)$ (the elements $u$ with $u$ and $u^{-1}$ integral)
this is Borel's theorem \cite[Thm.~5.1]{Borel1963} (see also \cite[\S5.1]{PlatonovRapinchuk1994}), which
the development proves (\lref{GlobalTheory.T.classNumber_finite_integralLevel_closedSubgroupGL}). For
general $U$: $U_0$ is compact, being closed (the locus is defined by polynomial equations) in the compact set
of $u$ with $u,u^{-1}$ integral; and $U$ is open, so $U_0$ is covered by finitely many cosets $tU$, $t\in T_1$. If
$x=\gamma su_0$ with $s\in T_0$, $u_0\in U_0$, and $u_0=tu$, then $x=\gamma(st)u$; take $T=T_0T_1$.
\end{proof}

\begin{corollary}\label{pf:geo:defclassfin}
Let $G$ be a unitary similitude group for $p$ with exactly one archimedean place of signature $(1,n-1)$,
$K\subseteq G(\A_f^p)$ a level, $K'\subseteq G_p$ a subgroup commensurable with $K_p$, and $\beta$ a definite form
(\cref{pf:ks1:def-forms,pf:ks1:def-pairs}) with its embedding $\iota_\beta$ of the definite group $\Gamma_\beta$ into
$G(\A_f^p)\times G_p$. Then the double coset space $\iota_\beta(\Gamma_\beta)\backslash(G(\A_f^p)\times G_p)/(K\times K')$ is
finite.
\formallabel{GlobalTheory.T.finite_doubleCoset_definiteHermitianForm_level}\formallabel{GlobalTheory.DefClassFin.ccf_holds}\formallabel{GlobalTheory.DefClassFin.cover_of_relIndex}\formallabel{GlobalTheory.DefClassFin.finite_of_cover}
\end{corollary}

\begin{proof}
The group $\{(\gamma,s)\in C^\times\times\Q^\times:\gamma\beta\gamma^*=s\beta\}$ (and more generally its intersection with the
centralizer of any element of $C$) is the group of rational points of a group locus over $\Q$
(\lref{GlobalTheory.T.exists_polynomialLocus_similitudeCentralizerTorus}: these similitude centralizers are
cut out in some $\GL_N$ by polynomial equations, functorially in the coefficient ring). Hence
\cref{pf:geo:classnumber} gives, for every open subgroup of its points in
$(\A_f^p\otimes C)^\times\times(\A_f^p)^\times$ times $(\Q_p\otimes C)^\times\times\Q_p^\times$, a finite set of representatives
$\gamma s u$ of the kind stated there. Applied with the open subgroup determined by $K\times K_p$, and transported through
the identification of $\beta$ with the standard form by the trivializing elements of the definite form, this covers
$G(\A_f^p)\times G_p$ by finitely many double cosets $\iota_\beta(\Gamma_\beta)\,x\,(K\times K_p)$
(\lref{GlobalTheory.DefClassFin.cover_KHT}). Since $K'$ and $K_p$ are commensurable, each double coset for
$K\times K_p$ is a union of finitely many double cosets for $K\times(K'\cap K_p)$, each of which is contained in a double
coset for $K\times K'$; and a set covered by finitely many double cosets has finitely many double cosets.
\end{proof}

% ===========================================================================
\subsection{Existence of plans}\label{sec:pf:geo:plans}

\begin{theorem}\label{pf:geo:planex}
$\sfname{PlanEx}$ holds.
\formallabel{UpperTheory.Gal.Int.Geo.Pl.planEx_holds}\formallabel{UpperTheory.Gal.Int.Geo.Pl.exists_unitaryGroupDatum_of_vector_exponent_spectral}
\end{theorem}

The proof assembles a plan from fifteen components. We first describe the objects they involve,
then state the components, prove the \emph{master lemma} (\cref{pf:geo:master}), which constructs a plan
with prescribed properties at two finite sets of places, and deduce \cref{pf:geo:planex}.

\subsubsection{Arithmetic data, setups, frames and vector data}

In this subsection $E$ is a CM field and $n\ge1$. Let $\geoGU_n(\A_f)$ be the group of pairs
$(g,\nu)\in\GL_n(\A_{E,f})\times\A_{\Q,f}^\times$ with ${}^t\bar gg=\nu$ (finite adelic points of the unitary
similitude group of the standard hermitian form on $E^n$), $\mathrm{comp}_w(g)\in\GL_n(E_w)$ the component of $g$ at a
place $w$, and, for a split place $w$ and $t\in\GL_n(E_w)$, $\mathrm{loc}_w(t)\in\geoGU_n(\A_f)$ the element of the unitary
group with component $t$ at $w$ (and the component at $w^c$ determined by the hermitian condition) and trivial
components elsewhere.

\begin{definition}[Arithmetic data]\label{pf:geo:arith}
An \emph{arithmetic datum} $A$ for $E$ consists of $\omega\in\OO_E$ and $\mathrm{tr},\mathrm{nm}\in\Z$ with
$\omega+\bar\omega=\mathrm{tr}$, $\omega\bar\omega=\mathrm{nm}$ and $\mathrm{tr}^2-4\mathrm{nm}<0$ (so $K=\Q(\omega)$ is imaginary quadratic); an
embedding $\tau_0\colon E\to\C$; a totally real $c\in\OO_E$ with $\tau_0(c)<0$ and $\tau(c)>0$ at the other archimedean
places; and $\delta\in E$ with $\bar\delta=-\delta$, $\mathrm{Tr}_{E/\Q}(\delta\OO_E)\subseteq\Z$ and $\Im\tau(\delta)<0$ for every $\tau$ with
$\tau(\omega)=\tau_0(\omega)$. A prime $q$ \emph{splits for $A$} if $q\nmid\mathrm{tr}^2-4\mathrm{nm}$ and $q\mid r^2-\mathrm{tr}\,r+\mathrm{nm}$ for some
$r\in\Z$; it is \emph{good for $A$} if moreover a perfectness condition holds at $q$ (there is $e$ prime to $q$ such
that $e$ times every $\Z$-linear form on $\OO_E$ is of the form $y\mapsto\mathrm{Tr}(\delta xy)$, and $c$ divides $e$). An element
$x\in E$ \emph{generates} $A$ if $\Q(x)=\Q(\omega)$.
\formal{UpperTheory.Gal.W2.Seed.Arith}\formal{UpperTheory.Gal.W2.Seed.Arith.SplitK}\formal{UpperTheory.Gal.W2.Seed.Arith.Good}\formal{UpperTheory.Gal.Int.Geo.Pl.Gen}
\end{definition}

\begin{definition}[Setups]\label{pf:geo:setup}
Let $A$ be an arithmetic datum, $n\ge1$ and $p$ a prime good for $A$. A \emph{setup} $S$ for $(A,n,p)$ consists of auxiliary
conditions on $A$ (at $n$ and at $p$) which make the following constructions possible, and of a link $L$ between the Harris--Taylor
datum $D_A$ and the unitary similitude group $G_A$ built from $A$ (\cref{pf:ks2:def:datum}), with $F_{D_A}=E$, $L$ inducing the
identity of $E$, and $\Lambda=\OO_E^n$ on which matrices with entries in $\OO_E$ act by multiplication. We write $D,G$ for $D_A,G_A$ and
$K_N\subseteq G(\A_f^p)$ for the principal congruence subgroups.
\formal{UpperTheory.Gal.W2.Seed.Setup}
\end{definition}

\begin{definition}[Frames]\label{pf:geo:frame}
A \emph{frame} for a setup $S$ consists of: a surjective homomorphism $j\colon\geoGU_n(\A_f)\to G(\A_f^p)$ whose kernel
consists of the $g$ with $\mathrm{comp}_w(g)=1$ for all $w\nmid p$; elements $a_w\in\GL_n(E_w)$, integral at the places above
primes good for $A$ and at almost all places; additive identifications of $\A_f^p\otimes V_\Q$ with the vectors of the
$E_w^n$, $w\nmid p$, which identify $\hat\Lambda^p$ with the integral vectors and under which $j(g)$ acts at $w$ by
$a_w\,\mathrm{comp}_w(g)\,a_w^{-1}$; compatibility with similitude factors; a definite form $\beta$ for $G$ whose definite
group is identified with the rational points of $\geoGU_n$, compatibly with $j$ and with the integrality at $p$; and the
involution $g\mapsto{}^tg^{-1}$ with inverse similitude factor.
\formal{UpperTheory.Gal.Int.Geo.Pl.Frame}
\end{definition}

Thus a frame identifies $G(\A_f^p)$ with the finite adelic points away from $p$ of the definite unitary similitude
group of $E^n$: Hecke operators and level groups of $\geoUa$ can be transported to $G$, and the definite inner forms of
$G$ become the definite group itself.

\begin{definition}[Vector data]\label{pf:geo:vec}
Let $x\in E$, $U$ an automorphic space on $\Uni_n(\A_{E^+})$, $w_0$ a split place and $\mathrm{Up}$ a predicate on places. A
\emph{vector datum} $V$ for $(x,U,w_0,\mathrm{Up})$ consists of: a weight $W$ (a dominant weight at each archimedean place)
and integers $m_v\ge0$, $j_z$ determined by $W$ and the $m_v$, with the associated weight
$\lambda_V\colon\Hom(E,\C)\to\N^n$ ($\lambda_V(\tau)$ is constant equal to $m_v$ at the chosen embedding $\tau$ of each archimedean place
$v$, and equal to $W_v+m_v$ at the other); an open compact torsion-free subgroup $J\subseteq\GL_n(E_{w_0})$ with a character
$\chi_J$ (open kernel) such that every smooth irreducible representation containing $\chi_J$ is supercuspidal; a finite
set $S$ of places containing $w_0$ and compact open subgroups $K_w\subseteq\GL_n(E_w)$ ($w\in S$), with $K_{w_0}=J$ and
$K_w\subseteq\GL_n(\OO_w)$ otherwise, such that every split $w\in S$ satisfies $\mathrm{Up}(w)$ and $w^c\notin S$; and a nonzero
function $\Phi$ on $\geoGU_n(\A_f)$ with values in the dual of the representation of weight $W$ (an algebraic modular form)
which is equivariant on the left under the rational points, transforms on the right under the level group
$\mathrm{lev}(S,(K_w))$ by $\chi_J$ at $w_0$, and is a spherical eigenvector with Satake multiset $B$ at every split $w$ with
$w,w^c\notin S$ and every Satake multiset $B$ of $U$ at $w$; every split place apart from $w_0$ at which $U$ has a Satake
multiset lies, with its conjugate, outside $S$; and a Hecke character $\kappa_V$, algebraic of infinity type
$m_V$, $m_V(\phi)=j_z-\sum_{v:\tau_v(x)=-\phi(x)}\sum_iW_{v,i}$ ($\tau_v$ the chosen embedding of the archimedean place $v$), unramified at
every $w$ such that $u^c\not\sim_xw$ for all $u\in S$, by which $\Phi$ transforms under the central elements attached to such $w$.
\formal{UpperTheory.Gal.Int.Geo.Pl.Vec}\formal{UpperTheory.Gal.Int.Geo.Pl.Vec.lam}
\end{definition}

For a vector datum $V$ and a place $w$: $w$ is \emph{free for $V$} if $w$ is split, $w\notin S$, $w^c\notin S$,
$u^c\not\sim_xw$ for all $u\in S$, and $w^c\not\sim_xw$; and $V$ is \emph{off $p$} if no place of $S$ contains $p$ and there is
$g_0$ with $\Phi(g_0)\ne0$ and $\mathrm{comp}_w(g_0)\in\GL_n(\OO_w)$ for all $w\mid p$.

\begin{definition}[Sites and seen Hecke data]\label{pf:geo:site}
Let $S$ be a setup, $F$ a frame, $V$ a vector datum, $N$ a level (an integer $N\ge3$ prime to $p$) and $w$ a place.
\begin{enumerate}[label=(\alph*),nosep]
\item A \emph{site} $\sigma$ at $w$ of level $N$ consists of a prime $q_\sigma$ good for $A$, $r_\sigma\in\Z$, an ideal $\mathfrak
  a_\sigma$ of $\OO_E$ and a subgroup $K_\sigma\subseteq G(\A_f^p)$ such that $(q_\sigma,r_\sigma,w)$ is admissible
  (\cref{pf:ks2:def:goodgrp}), $(N',\mathfrak a_\sigma)$ is an auxiliary level for it, where $N'$ is $N$ with its
  $q_\sigma$-part removed, $K_\sigma$ is the level group of the ideal $N'\mathfrak a_\sigma$, and $N\in N'\mathfrak a_\sigma$. Sites $\sigma$, $\sigma'$
  are \emph{compatible} if $q_\sigma=q_{\sigma'}$ implies $r_\sigma=r_{\sigma'}$ and $\mathfrak a_\sigma=\mathfrak a_{\sigma'}$.
\item A Hecke datum $g=(K_g,h_g,M_g)$ of level $N$ is a \emph{deck member} if $K_g=K_N$, $M_g=N$ and $h_g\in K_1$; it is a
  \emph{$T_i$-member at the site $\sigma$} if $i\le n$, $K_g=K_\sigma$, $M_g=N(\Nm w)^i$ and $h_g=j(\mathrm{loc}_w(t)z^i)$ with
  $t=\mathrm{diag}(\varpi,\dots,\varpi,1,\dots,1)$ ($i$ entries $\varpi$, $\varpi$ a uniformizer at $w$) and $z\in\geoGU_n(\A_f)$
  central with $\mathrm{comp}_u(z)=\Nm w$ if $u^c\sim_xw$ and $\mathrm{comp}_u(z)=1$ otherwise. Its \emph{index} is
  $\mathrm{idx}(g)=[K_{M_g}:K_g\cap h_g^{-1}K_gh_g]$ (relative index).
\item An element $k\in G(\A_f^p)$ is \emph{stable} if $k=j(k')$ with ${}^tk'^{-1}$ in the level group $\mathrm{lev}(S,(K_w))$ of $V$
  and $\chi_J$ trivial on its component at $w_0$.
\item For a prime $r_0$, a pair $(g,\theta)$ of a Hecke datum of level $N$ and $\theta\in\C$ is \emph{seen} if either $g$ is a deck
  member with $h_g$ stable, $h_g\in K_{r_0}$ and $\theta=1$; or there are a place $w$ free for $V$, a site $\sigma$ at $w$ with
  $q_\sigma\ne r_0$ and all elements of $K_\sigma$ stable, $i\le n$ and a Satake multiset $B$ of $U$ at $w$ such that $g$ is a
  $T_i$-member at $\sigma$ and
  \[
    \theta=\mathrm{idx}(g)\cdot(\Nm w)^{i(n-i)/2}e_i(B)\cdot\kappa_V(w)^{-i},
  \]
  $e_i$ the $i$-th elementary symmetric function.
\end{enumerate}
\formal{UpperTheory.Gal.Int.Geo.Pl.Site}\formal{UpperTheory.Gal.Int.Geo.Pl.Site.Compat}\formal{UpperTheory.Gal.Int.Geo.Pl.IsDeckM}\formal{UpperTheory.Gal.Int.Geo.Pl.IsTM}\formal{UpperTheory.Gal.Int.Geo.Pl.idx}\formal{UpperTheory.Gal.Int.Geo.Pl.Stab}\formal{UpperTheory.Gal.Int.Geo.Pl.Seen}
\end{definition}

The number $(\Nm w)^{i(n-i)/2}e_i(B)$ is the eigenvalue of the Hecke operator $T_i$ on a spherical vector with Satake
multiset $B$; the factor $\mathrm{idx}(g)$ and the power of $\kappa_V(w)$ account for the normalisation of the Hecke data and
for the twist by $\kappa_V$.

\subsubsection{The components}

The proof uses the following fifteen statements; they are the components of the bundle $\sfname{All}$ of the
development. Four of them are proved below (\cref{pf:geo:easycomp}); the others are proved in the formal development
and are cited.
\begin{description}[style=unboxed,leftmargin=1em,nosep]
\item[$\sfname{DatEx}$] (\lref{UpperTheory.Gal.Int.Geo.Pl.datEx_holds}). Let $n$ be odd, $\sfname{FieldOK}(E,\ell,d)$,
  $x^2=-d$, and $Q$ a finite set of primes $q$ such that $-d$ is a square in $\Q_q$. Then there is an arithmetic datum $A$
  generated by $x$ such that every $q\in Q$ is good, there is $e>0$ such that every prime $q\nmid e$ for which $-d$ is a
  square in $\Q_q$ is good, and for every finite set $X$ of primes there are a prime $p\notin X$ and a setup for $(A,n,p)$
  satisfying the standing conditions $\mathrm{Sc}$.
\item[$\sfname{Sides}$] (\cref{pf:geo:easycomp}). The splitting behaviour of good primes, see below.
\item[$\sfname{VecEx}$] (\lref{UpperTheory.Gal.Int.Geo.Pl.vecEx_holds}). Let $n,d>0$, $x^2=-d$, $U$ automorphic,
  $w_0$ split, $U$ with a supercuspidal type of $\ell$-power order at $w_0$, and $\mathrm{Up}$ a predicate such that
  $\mathrm{Up}(w)$ or $\mathrm{Up}(w^c)$ for every split $w$, and $\mathrm{Up}(w_0)$. Then a vector datum for $(x,U,w_0,\mathrm{Up})$
  exists.
\item[$\sfname{IntCof}$] (\cref{pf:geo:easycomp}). Every $g\in\geoGU_n(\A_f)$ has integral components at almost all places.
\item[$\sfname{FrameEx}$] (\lref{UpperTheory.Gal.Int.Geo.Pl.frameEx_holds}). For $n$ odd every setup has a frame.
\item[$\sfname{SpecTop}$] (\lref{UpperTheory.Gal.Int.Geo.Pl.specTop_holds}). Let $n$ be odd, $S$ a setup with
  $\mathrm{Sc}$, $F$ a frame, $x$ a generator, $V$ a vector datum off $p$, and assume some good prime lies in $w_0$.
  For every finite set $X$ of primes there is a prime $r_0\notin X$, $r_0\ne p$, such that: for every level $N$ with $r_0\mid N$
  and all elements of $K_N$ stable, and every family $(g_j,\theta_j)_{j\in J}$ of seen pairs, there is a function $f$ in the Hecke
  algebra of $K_N$, supported on elements with integral inverse, with the spectral condition
  $\mathcal S(L,\lambda_V,K_N,(f_{g_j})_j,(\theta_j)_j,f)$ of \cref{pf:ks1:def-spec}.
\item[$\sfname{SiteEx}$] (\lref{UpperTheory.Gal.Int.Geo.Pl.siteEx_holds}). Let $x$ be a generator, $V$ off $p$, $Q$ a
  finite set of places and $q\colon Q\to\{\text{primes}\}$ with: $q(w)\ne p$ good, $q(w)\in w$, $w$ free for $V$, and $w\sim_xw'$
  whenever $q(w)=q(w')$. For every $c>0$ prime to $p$ there are a level $N$ divisible by $c$ with all elements of $K_N$
  stable, and sites $\sigma_w$ at $w\in Q$ with $q(w)\mid N$, $q_{\sigma_w}=q(w)$, all elements of $K_{\sigma_w}$ stable, and
  $\sigma_w,\sigma_{w'}$ compatible for all $w,w'$.
\item[$\sfname{SiteTriv}$] (\lref{UpperTheory.Gal.Int.Geo.Pl.siteTriv_holds}). If $q\ne p$ is a good prime not
  dividing the level $N$ and $q\in w$, there is a site $\sigma$ at $w$ with $q_\sigma=q$ and $K_\sigma=K_N$.
\item[$\sfname{LevBasic}$] (\lref{UpperTheory.Gal.Int.Geo.Pl.levBasic_holds}). Among other elementary properties: level
  groups of ideals exist and are monotone in the ideal, $K_N$ is the level group of $(N)$, level groups are normalized by $K_1$, and if $h$
  is integral and $h\hat\Lambda^p\supseteq c\hat\Lambda^p$ with $c>0$, then $hK_{Nc}h^{-1}\subseteq K_N$.
\item[$\sfname{MemEx}$] (\lref{UpperTheory.Gal.Int.Geo.Pl.memEx_holds}). For a generator $x$ and a level $N$: every
  $k\in K_1$ is $h_g$ for a deck member $g$; and for every site $\sigma$ at a place $w$ and every $i\le n$ there is a $T_i$-member
  at $\sigma$.
\item[$\sfname{ExpOf}$] (\lref{UpperTheory.Gal.Int.Geo.Pl.expOf_holds}). For a vector datum $V$ and a generator $x$
  with $x^2=-d$, $d>0$: each $\lambda_V(\tau)$ is non-increasing, and for every $\tau\colon E\to\C$ there is $\sigma\in\Aut(\C)$ with
  $\sigma\tau=\tau_0$ such that for every weight $\lambda_U$ of $U$ at $\tau$ and every $j\in J(\lambda_U)$ there are a Hodge exponent vector
  $h$ for $(\tau_0,\varepsilon_D,\sigma\lambda_V)$ and an index $i$ with $h_i=j+m_V(\tau)$.
\item[$\sfname{GoodOf}$] (\lref{UpperTheory.Gal.Int.Geo.Pl.goodOf_holds}). Let $\sigma$ be a site at $v$ and $(g_j,\theta_j)_j$
  a family of Hecke data with values in a field satisfying the deck condition for $K_\sigma$ (\cref{pf:ks2:def:goodgrp}),
  every $g_j$ being a deck member or a $T_i$-member at a site $\sigma'$ at a place $w\ne v$ compatible with $\sigma$. Then the
  family is unramified at $v$: $\mathcal U(g,\theta,v)$.
\item[$\sfname{CongOf}$] (\lref{UpperTheory.Gal.Int.Geo.Pl.congOf_holds}). Let $\sigma$ be a site at $w$, $(g_j,\theta_j)_j$
  a family with the deck condition for $K_\sigma$, and $\vartheta\colon\N\to L$ such that for every $i\le n$ some $g_j$ is a
  $T_i$-member at $\sigma$ with $\theta_j=\mathrm{idx}(g_j)\vartheta_i$. Then $\mathcal H(g,\theta,w,\vartheta)$.
\item[$\sfname{RootsAlg}$, $\sfname{SatCard}$] (\cref{pf:geo:easycomp}, \cref{pf:geo:norms}).
\end{description}

\begin{lemma}[$\sfname{Sides}$, $\sfname{IntCof}$, $\sfname{RootsAlg}$]\label{pf:geo:easycomp}
\begin{enumerate}[label=(\alph*),nosep]
\item Let $A$ be an arithmetic datum, $x$ a generator of $A$, $q$ a prime which splits for $A$, and $w$ a place with $q\in w$.
  Then: $w^c\not\sim_xw$; $w$ is split; every place $u$ with $q\in u$ satisfies $u\sim_xw$ or $u^c\sim_xw$; $u\sim_xw$ implies
  $q\in u$; $u\sim_xu'$ implies $u^c\sim_xu'^c$; there is $r\in\Z$ with $q\mid r^2-\mathrm{tr}\,r+\mathrm{nm}$ and $\omega-r\in w$; and for every
  $r\in\Z$ with $\omega-r\in w$ and every place $u$, $u\sim_xw$ if and only if $q\in u$ and $\omega-r\in u$.
\item For every $g\in\geoGU_n(\A_f)$, $\mathrm{comp}_w(g)\in\GL_n(\OO_w)$ for all but finitely many $w$.
\item Let $q,n\in\N$, $q>0$, $B$ a multiset of $n$ nonzero complex numbers, $a,c\in\C$, and
  $\vartheta_i=a\,q^{i(n-i)/2}e_i(B)\,c^i$ for $i\le n$. Then $\vartheta_0=a$ and
  $\sum_{i=0}^n(-1)^iq^{\binom i2}\vartheta_iX^{n-i}=a\prod_{b\in B}(X-cq^{(n-1)/2}b)$.
\end{enumerate}
\formallabel{UpperTheory.Gal.Int.Geo.Pl.sides_holds}\formallabel{UpperTheory.Gal.Int.Geo.Pl.intCof_holds}\formallabel{UpperTheory.Gal.CertificateRingHomMap.rootsAlg_holds}
\end{lemma}

\begin{proof}
(a) Let $K=\Q(x)=\Q(\omega)$; then $\OO_E\cap K=\OO_K$, and $w\sim_xu$ means $w\cap\OO_K=u\cap\OO_K$. Since $q$ does not divide the
discriminant $\mathrm{tr}^2-4\mathrm{nm}$ of $\Z[\omega]$, it does not divide $[\OO_K:\Z[\omega]]$, so $\OO_K/q\OO_K\cong\Z[\omega]/(q)\cong\F_q[X]/(X^2-\mathrm{tr}X+\mathrm{nm})$;
the polynomial has two distinct roots modulo $q$, hence $q\OO_K=\mathfrak q\bar{\mathfrak q}$ with $\mathfrak q\ne\bar{\mathfrak q}$, $\mathfrak q=(q,\omega-r)$
and $\bar{\mathfrak q}=(q,\omega-r')$ for the two roots $r,r'$. The place $w$ lies over one of them, say $w\cap\OO_K=\mathfrak q$, and $w^c$ over
$\bar{\mathfrak q}$; so $w^c\not\sim_xw$, in particular $w^c\ne w$, and $w$ is split. Every $u\ni q$ lies over $\mathfrak q$ or $\bar{\mathfrak q}$, i.e.\
$u\sim_xw$ or $u^c\sim_xw$. If $u\sim_xw$ then $q\in w\cap\OO_K\subseteq u$. Conjugation preserves $\sim_x$. Since $\OO_K/\mathfrak q=\F_q$, there is
$r\in\Z$ with $\omega-r\in\mathfrak q\subseteq w$, and then $r^2-\mathrm{tr}\,r+\mathrm{nm}\in\mathfrak q\cap\Z=q\Z$. Finally, if $\omega-r\in w$, then $(q,\omega-r)\OO_K$ is
the maximal ideal $w\cap\OO_K$ (it is contained in it and is maximal by the description of $\OO_K/q\OO_K$); so $q,\omega-r\in u$ implies
$u\cap\OO_K\supseteq w\cap\OO_K$, hence equality, and conversely.

(b) The matrix $\mathrm{comp}_w(g)$ and its inverse are the $w$-components of elements of $\GL_n$ of the finite adeles, which are
integral at almost all places.

(c) $\vartheta_0=a$ since $e_0=1$. For each $i$, $q^{\binom i2}q^{i(n-i)/2}=q^{i(n-1)/2}$, so the left side is
$a\sum_i(-1)^ie_i(B)(cq^{(n-1)/2})^iX^{n-i}=a\prod_{b\in B}(X-cq^{(n-1)/2}b)$.
\end{proof}

\subsubsection{The master lemma}

\begin{proposition}[Master lemma]\label{pf:geo:master}
Assume the fifteen components. Let $E$ be a CM field, $\ell$ a prime, $\iota\colon\Qlbar\to\C$, $a$ odd, $\sfname{FieldOK}(E,\ell,d)$,
$x^2=-d$, $w_0$ a place with $\ell\in w_0$, and $U$ automorphic with a supercuspidal type of $\ell$-power order at $w_0$. Let
$H,C$ be disjoint finite sets of places such that every $w\in H\cup C$ is $d$-usable, apart from $w_0$, $U$ has a Satake multiset
at $w$, $w\sim_xw_0$ if $p_w=\ell$, and $w\sim_xw'$ for all $w'\in H\cup C$ with $p_{w'}=p_w$. Then there is a plan $P$ for $(\iota,d,U)$
with $\mathrm{HTFor}(P,v)$ for every $v\in H$ and $\mathrm{CongAll}(P,w)$ for every $w\in C$.
\formallabel{UpperTheory.Gal.Int.Geo.Pl.master}
\end{proposition}

\begin{proof}
\emph{Step 1: arithmetic datum.} For a $d$-usable place $w$, $-d$ is a square in $\Q_{p_w}$; and $-d$ is a square in $\Ql$ by
$\sfname{FieldOK}$. Apply $\sfname{DatEx}$ with $Q=\{\ell\}\cup\{p_w:w\in H\cup C\}$: we obtain an arithmetic datum $A$ generated by
$x$, with $\ell$ and all $p_w$ ($w\in H\cup C$) good, an integer $e>0$ as in $\sfname{DatEx}$, and the supply of auxiliary primes with
setups. By \cref{pf:geo:easycomp}(a), applied to $w_0$ and $q=\ell$, and to $w$ and $q=p_w$ for $w\in H\cup C$, all these places
are split, $w_0^c\not\sim_xw_0$ and $w^c\not\sim_xw$.

\emph{Step 2: the side predicate.} Put $\mathrm{Up}(u)$: $u^c\not\sim_xw$ for all $w\in H\cup C$. We claim that for every split $u$,
$\mathrm{Up}(u)$ or $\mathrm{Up}(u^c)$. Otherwise there are $w,w'\in H\cup C$ with $u^c\sim_xw$ and $u\sim_xw'$; then
$p_w=p_{u^c}=p_u=p_{w'}$, so $w\sim_xw'$ by the coherence hypothesis and $u\sim_xw$; by \cref{pf:geo:easycomp}(a) (fifth clause) $u^c\sim_xw^c$,
so $w^c\sim_xu^c\sim_xw$, a contradiction. Moreover $\mathrm{Up}(w_0)$: if $w_0^c\sim_xw$ for some $w\in H\cup C$, then $p_w=\ell$, so
$w\sim_xw_0$ and $w_0^c\sim_xw_0$, a contradiction.

\emph{Step 3: vector datum.} By $\sfname{VecEx}$ (with $n=a$, $w_0$ split) there is a vector datum $V$ for $(x,U,w_0,\mathrm{Up})$. Every
$w\in H\cup C$ is free for $V$: it is split; $w,w^c\notin S$ because $w$ is split, apart from $w_0$ and $U$ has a Satake multiset at $w$;
if $u\in S$ and $u^c\sim_xw$, then $p_w\in u$, so $u$ is split by \cref{pf:geo:easycomp}(a), hence $\mathrm{Up}(u)$ by the definition of
vector data, contradicting $u^c\sim_xw$; and $w^c\not\sim_xw$.

\emph{Step 4: auxiliary prime $p$, setup and frame.} Choose $g_0$ with $\Phi(g_0)\ne0$ and let $F_0$ be the finite set of places at
which $\mathrm{comp}(g_0)$ is not integral (\cref{pf:geo:easycomp}(b)). By $\sfname{DatEx}$ there are a prime $p$ with
$p\notin\{p_w:w\in H\cup C\}\cup\{p_u:u\in S\cup F_0\}$ and a setup $S$ for $(A,a,p)$ with $\mathrm{Sc}(D,G)$. Then $V$ is off $p$. By
$\sfname{FrameEx}$ there is a frame $F$ for $S$.

\emph{Step 5: the spectral prime and the level.} Apply $\sfname{SpecTop}$ (a good prime, namely $\ell$, lies in $w_0$) with
$X=\{p_w:w\in H\cup C\}$: we get a prime $r_0\notin X$, $r_0\ne p$. By $\sfname{SiteEx}$ with $Q=H\cup C$, $q(w)=p_w$ (the hypotheses
were verified in Steps 1--4) and $c=r_0$, there are a level $N$ with $r_0\mid N$, all elements of $K_N$ stable, and pairwise
compatible sites $\sigma_w$ ($w\in H\cup C$) with $q_{\sigma_w}=p_w\mid N$ and all elements of $K_{\sigma_w}$ stable. Note $p_w\ne r_0$.

\emph{Step 6: the Hecke data.} Let $J$ be the set of seen pairs $(g,\theta)$ (with respect to $F,V,r_0,N$) such that $g$ is a deck member,
or $g$ is a $T_i$-member at a site $\sigma'$ at a place $w\notin H$ such that every $\sigma_v$, $v\in H$, is compatible with $\sigma'$. Put
$g_j=g$ and $\theta_j=\iota^{-1}(\theta)$ for $j=(g,\theta)$. By $\sfname{SpecTop}$ there is $f$ in the Hecke algebra of $K_N$, supported on
elements with integral inverse, with $\mathcal S(L,\lambda_V,K_N,(f_{g_j}),(\iota\theta_j),f)$. By $\sfname{ExpOf}$ each $\lambda_V(\tau)$ is
non-increasing.

\emph{Step 7: the deck condition.} Let $\sigma'$ be a site at some $w$ with $q_{\sigma'}\ne r_0$ and all elements of $K_{\sigma'}$ stable. Then
$(g_j,\theta_j)$ satisfies the deck condition for $K_{\sigma'}$. Indeed, let $k\in K_{\sigma'}$; $k\in K_1$, and $k\in K_{r_0}$ by the monotonicity
of level groups ($\sfname{LevBasic}$), because $r_0$ divides $N$ but not $q_{\sigma'}$, so the ideal $N'\mathfrak a_{\sigma'}$ is contained in $(r_0)$. By
$\sfname{MemEx}$ there is a deck member $g$ with $h_g=k$, and $(g,1)$ is seen ($k$ stable, $k\in K_{r_0}$), so $(g,1)\in J$ with $\theta=1$.

\emph{Step 8: congruence data.} Let $w$ be free for $V$, $\sigma'$ a site at $w$ with $q_{\sigma'}\ne r_0$ and all elements of $K_{\sigma'}$
stable, $w\notin H$, every $\sigma_v$ ($v\in H$) compatible with $\sigma'$, and $B$ a Satake multiset of $U$ at $w$. We show
$\sfname{CongP}(w,B)$ for the data being assembled (with $e=\mathrm{id}$, $\mathfrak w=w$). The character $\kappa_V$ is unramified at $w$ since $w$
is free. Put $c=\kappa_V(w)^{-1}$ and $\vartheta_i=\iota^{-1}\bigl((\Nm w)^{i(a-i)/2}e_i(B)c^i\bigr)$. For each $i\le a$, $\sfname{MemEx}$ gives a
$T_i$-member $g$ at $\sigma'$, and $(g,\mathrm{idx}(g)(\Nm w)^{i(a-i)/2}e_i(B)c^i)$ is seen and lies in $J$; its value is $\mathrm{idx}(g)\vartheta_i$.
Together with Step~7, $\sfname{CongOf}$ gives $\mathcal H(g,\theta,w,\vartheta)$. By \cref{pf:geo:easycomp}(c) with $q=\Nm w$, $n=a$, $a=1$
(using $|B|=a$, (S12)), $\iota(\vartheta_0)=1\ne0$ and
$\iota\bigl(\sum_i(-1)^i(\Nm w)^{\binom i2}\vartheta_iX^{a-i}\bigr)=\prod_{b\in B}(X-c(\Nm w)^{(a-1)/2}b)$, i.e.\ $\kappa_V(w)^{-1}(\Nm w)^{(a-1)/2}B$ is an
$\iota$-root multiset for $(w,a,\vartheta)$.

\emph{Step 9: the plan.} Let $P$ consist of $p,D,G,L$ (with $\mathrm{Sc}$), $e=\mathrm{id}_E$, $N$, $(g_j)$, $\theta$, $\kappa=\kappa_V$ with $m=m_V$,
$\lambda=\lambda_V$, $f$, and the finite set $X'$ of places $w$ such that $p_w$ divides $Npe$ or equals $p_u$ for some $u\in S$. For $X'$: let
$w\notin X'$ be $d$-usable with $\ell\notin w$ and $B$ a Satake multiset at $w$. Then $p_w\nmid Npe$, so $p_w$ is good by $\sfname{DatEx}$, and
$\sfname{SiteTriv}$ gives a site $\sigma'$ at $w$ with $q_{\sigma'}=p_w$ and $K_{\sigma'}=K_N$. The place $w$ is free for $V$ (split; $w,w^c\notin S$ and
no $u\in S$ has $u^c\sim_xw$, since $p_u\ne p_w$; $w^c\not\sim_xw$ by \cref{pf:geo:easycomp}(a)); $q_{\sigma'}=p_w\ne r_0$ since $r_0\mid N$; the
elements of $K_N$ are stable; $w\notin H$ since $p_v\mid N$ for $v\in H$; and compatibility is vacuous since $q_{\sigma_v}=p_v\ne p_w$. So Step~8
gives $\sfname{CongP}(w,B)$, and $P$ is a plan.

\emph{Step 10: the places of $H$ and $C$.} Let $v\in H$. Then $\kappa_V$ is unramified at $v$ ($v$ is free). The site $\sigma_v$ has
$q_{\sigma_v}=p_v\ne r_0$ and stable elements, so Step~7 gives the deck condition for $K_{\sigma_v}$; every $g_j$ is a deck member or a $T_i$-member at
a site at a place $w\notin H$ (so $w\ne v$) compatible with $\sigma_v$. By $\sfname{GoodOf}$, $\mathcal U(g,\theta,v)$. The exponent condition at $v$
follows from $\sfname{ExpOf}$ (the distinguished embedding of $D_A$ is $\tau_0$, and $i_0$ is the unique index). So $\mathrm{HTFor}(P,v)$. Let
$w\in C$; then $w\notin H$, $\sigma_w$ is compatible with every $\sigma_v$, $q_{\sigma_w}=p_w\ne r_0$ and its elements are stable, so Step~8 gives
$\sfname{CongP}(w,B)$ for every Satake multiset $B$ at $w$, i.e.\ $\mathrm{CongAll}(P,w)$.
\end{proof}

\begin{theorem}\label{pf:geo:planexof}
The fifteen components imply $\sfname{PlanEx}$.
\formallabel{UpperTheory.Gal.Int.Geo.Pl.planEx_of}
\end{theorem}

\begin{proof}
Let the data of $\sfname{PlanEx}$ be given. Every place $w$ with $\ell\in w$ is $d$-usable: it is split, since non-split places do not contain
$\ell$ ($\sfname{FieldOK}$), and $-d$ is a square in $\Ql$. Also $v'\sim_xw_0$. We apply the master lemma four times.
(i) $H=\{v'\}$, $C=\emptyset$: the conditions on $v'$ hold ($v'$ is usable, apart from $w_0$, has a Satake multiset, and $v'\sim_xw_0$); this gives a
plan with $\mathrm{HTFor}(P,v')$. (ii) For $w$ usable with $\ell\notin w$ and a Satake multiset at $w$: $H=\{v'\}$, $C=\{w\}$ ($w\ne v'$ as
$\ell\in v'$). The place $w$ is apart from $w_0$: otherwise $\ell$, which lies in the place of $E^+$ below $w_0$, would lie in $w$. Since $p_w\ne\ell=p_{v'}$,
the coherence conditions are empty. (iii) For $w$ as in (G3): $H=\{v',w\}$, $C=\emptyset$; the coherence conditions are $w\sim_xv'$ (given)
and $w\sim_xw_0$, which follows from $w\sim_xv'\sim_xw_0$. This gives $\mathrm{HTFor}$ at $v'$ and $w$. (iv) For the same $w$: $H=\{v'\}$, $C=\{w\}$
($w\ne v'$ since they are apart), giving $\mathrm{HTFor}(P,v')$ and $\mathrm{CongAll}(P,w)$.
\end{proof}

\begin{proof}[Proof of \cref{pf:geo:planex}]
Apply \cref{pf:geo:planexof}: $\sfname{Sides}$, $\sfname{IntCof}$ and $\sfname{RootsAlg}$ hold by \cref{pf:geo:easycomp}, $\sfname{SatCard}$
is (S12) (\cref{pf:geo:norms}), and the other eleven components hold by the results cited in the list above.
\end{proof}

\subsubsection{The spectral component}\label{sec:pf:geo:spec}

The component $\sfname{SpecTop}$ carries most of the weight of the construction of plans; it is where the Hecke eigensystem of
$U$ is shown to occur in the algebraic modular forms of the definite inner forms of $G$, detected by traces of good test
functions. Its formal proof (\lref{UpperTheory.Gal.Int.Geo.Pl.specTop_holds}, cited) has the following structure.
\begin{enumerate}[label=(\roman*),nosep]
\item \emph{Auxiliary prime} (\lref{UpperTheory.Gal.Spc.auxiliary_prime_holds}). Given $X$, a prime $r_0\notin X$, $r_0\ne p$, is chosen
  larger than all primes in $X$, than $p$, than the discriminant of $\omega$ and the norms of the places of $S$, and with a root of
  $X^2-\mathrm{tr}X+\mathrm{nm}$ modulo $r_0$; it is good for $A$, prime to the places of $S$, and comes with a set of places above it.
\item \emph{Type datum} (\lref{UpperTheory.Gal.Spc.tyEx}). The supercuspidal type $(J,\chi_J)$ at $w_0$ is transported through the
  frame to a type function $e$ in the Hecke algebra of a level $K_0$; $e$ is idempotent, and it satisfies the cuspidality condition
  used in (v) at the place $w_0$ (\lref{UpperTheory.Gal.Spc.cuspIrr_kFinite_rep_eq_centralCuspForm}), because every smooth
  irreducible representation containing the type is supercuspidal.
\item \emph{Regularizing function} (\lref{UpperTheory.Gal.Spc.auxFn}). At the places above $r_0$ a function $u$ with a convolution
  inverse $u'$ on the $e$-part is constructed, commuting with the Hecke functions of the seen Hecke data, such that for every $m\ge1$ the
  support of the $m$-fold power of $u$ on the $e$-part consists of elements whose components at an auxiliary place above $r_0$ have
  separable characteristic polynomial and which are regular (\cref{pf:ks1:def-htgroup}); this set of elements is stable under
  multiplication by the supports of the Hecke functions $f_{g_j}$.
\item \emph{Eigenvector and weight.} For the coefficient representation $\xi$ of the vector datum composed with the transpose-inverse,
  the $\xi$-valued forms of the $e$-part contain an eigenvector of the Hecke functions $f_{g_j}$ with eigenvalues $\theta_j$
  (\lref{UpperTheory.Gal.Spc.eigEx}), and $\xi$ has weight $\lambda_V$ in the sense of \cref{pf:ks1:def-spec}
  (\lref{UpperTheory.Gal.Int.Geo.Pl.wtOf_holds}).
\item \emph{Abstract spectral lemma} (\cref{pf:geo:specoff} below).
\end{enumerate}
The lemmas (i)--(iv) are proved in the formal development; their conversion is deferred. The abstract spectral lemma is proved
below, modulo lemmas on Hecke algebras and orbital integrals which are cited.

For a level $K$ of $G(\A_f^p)$, convolution makes the Hecke algebra of $K$ (bi-$K$-invariant compactly supported functions) a ring;
for a family $(\psi_j)_{j\in J}$ of Hecke functions and a finitely supported function $s$ on words in $J$ we write $s(\psi)*f$ as in
\cref{pf:ks1:def-spec}, and for Hecke functions $u,e$ we write $u^{*m}*e$ for the $m$-fold convolution of $u$ applied to $e$. Functions $e,u,u'$
satisfy the \emph{commutation relations} with $(\psi_j)$ if $e*e=e$, $e*\psi_j*e=\psi_j*e$, $e*u*e=u*e$, $\psi_j*\psi_{j'}*e=\psi_{j'}*\psi_j*e$,
$u*\psi_j*e=\psi_j*u*e$ and $u'*u*e=e$.

\begin{proposition}[Abstract spectral lemma]\label{pf:geo:specoff}
Let $G$ be a unitary similitude group for $p$ with one archimedean place of signature $(1,n-1)$, $K$ a level, $D,L$ a datum and a link,
$\lambda$ a weight, $\xi$ a finite-dimensional representation of $C^\times\times\Q^\times$ of weight $\lambda$, $(\psi_j)_{j\in J}$ Hecke functions of
level $K$, $\theta\colon J\to\C$, and $e,u,u'$ Hecke functions of level $K$ satisfying the commutation relations. Fix a place of $F$ at which a
cuspidality condition $\mathrm{C}$ on functions is given and an auxiliary place at which a separability condition $\mathrm R$ is given
(the characteristic polynomial of the component is separable), as in (ii) and (iii). Assume: the classes of definite forms are finite;
$e$ satisfies $\mathrm{C}$; for every $m\ge1$ there is a set $\mathcal A_m$ of elements satisfying $\mathrm R$ and regular, stable under
multiplication by the supports of the $\psi_j$, containing the support of $u^{*m}*e$; and some definite form class carries a nonzero
$\xi$-valued form $\Phi$ with $T(e)\Phi=\Phi$ and $T(\psi_j)\Phi=\theta_j\Phi$ for all $j$, $T$ denoting the Hecke operators on forms
(\cref{pf:ks1:def-dtrace}). Then there is $m\ge1$ with $\mathcal S(L,\lambda,K,(\psi_j),\theta,u^{*m}*e)$.
\formallabel{UpperTheory.Gal.Spc.specOfF}
\end{proposition}

\begin{lemma}\label{pf:geo:specalg}
Let $(W_q)_{q\in Q}$ be a finite family of finite-dimensional complex vector spaces, and for each $q$ let $T_q(j)$ ($j\in J$), $e_q,u_q,u'_q$ be
endomorphisms of $W_q$ satisfying the commutation relations (for composition). Assume that some $W_{q_0}$ contains $x_0\ne0$ with
$e_{q_0}x_0=x_0$ and $T_{q_0}(j)x_0=\theta_jx_0$ for all $j$. Then there are $m\ge1$ and a finitely supported function $a$ on words such that
for every $s$ with $s(\theta)\ne0$
\[
  \sum_{q\in Q}\tr\bigl((as)(T_q)\,u_q^m\,e_q\mid W_q\bigr)\ne0,
\]
where $as$ is the product of $a$ and $s$ (concatenation of words).
\formallabel{UpperTheory.Gal.Spc.specAlg}\formallabel{UpperTheory.Gal.CertificateRingHomMap.pTracePow_holds}
\end{lemma}

\begin{proof}
Let $W=\bigoplus_qW_q$ and let $t_j$, $v$, $E$, $v'$ be the block diagonal operators with blocks $T_q(j)e_q$, $u_qe_q$, $e_q$, $u'_q$. By the
commutation relations the $t_j$, $v$ and $E$ commute pairwise, $E$ is idempotent and $v'v=E$. By \cref{pf:geo:projone} there is $a$ with
$a(\theta)=1$, $P=a(t)$ idempotent and $(t_j-\theta_j)^NP=0$ for some $N$. The vector $x=(0,\dots,x_0,\dots,0)$ satisfies $t_jx=\theta_jx$ and
$Ex=x$, hence $Px=a(\theta)x=x$; so the idempotent $PE$ is nonzero. The operator $v$ commutes with $PE$ and is injective on $W'=PE(W)$,
since $v'v=E$ is the identity on $W'$. Hence $v|_{W'}$ is invertible, and some power sum $\tr((v|_{W'})^m)$ with $m\ge1$ is nonzero: otherwise
all power sums of the eigenvalues of $v|_{W'}$ would vanish, hence (Newton's identities, characteristic $0$) all elementary symmetric
functions, and $v|_{W'}$ would be nilpotent. Now let $s(\theta)\ne0$. On $W'$ each $t_j$ acts as $\theta_j$ plus a nilpotent operator commuting with
everything, so $s(t)$ acts on $W'$ as $s(\theta)$ plus a nilpotent operator commuting with $v$, and
$\tr(s(t)\,P\,v^mE)=s(\theta)\tr((v|_{W'})^m)\ne0$. Finally, by the commutation relations
$\sum_q\tr((as)(T_q)u_q^me_q)=\tr(a(t)s(t)v^mE)$, since $(T_q(j)e_q)$-words times $e_q$ equal $T_q(j)$-words times $e_q$ and
$(u_qe_q)^me_q=u_q^me_q$.
\end{proof}

\begin{proof}[Proof of \cref{pf:geo:specoff}]
The convolutions $u^{*m}*e$ and $s(\psi)*u^{*m}*e$ are Hecke functions of level $K$, and $\psi_j*\psi_{j'}*\psi_w*u^{*m}*e=\psi_{j'}*\psi_j*\psi_w*u^{*m}*e$
for all words $w$ (\lref{UpperTheory.Gal.Spc.words}); this is the commutativity requirement of $\mathcal S$. For each class $q$ of
definite forms (a finite set) let $W_q$ be the space of $\xi$-valued forms of the class (finite-dimensional by \cref{pf:geo:defclassfin}); the
Hecke operators give an algebra homomorphism $f\mapsto T_q(f)$ from the Hecke algebra of $K$ to $\End(W_q)$
(\lref{UpperTheory.Gal.Spc.opAlg}), and by definition $\mathrm{Tr}^{\mathrm{def}}(\xi;K,f;K_p,\mathbf 1_{K_p})=\sum_q\tr(T_q(f)\mid W_q)$. The
eigenvector of the hypothesis can be moved to the chosen representative of its class (\lref{UpperTheory.Gal.Spc.clsMove}). The
images $T_q(\psi_j),T_q(e),T_q(u),T_q(u')$ satisfy the commutation relations, so \cref{pf:geo:specalg} gives $m\ge1$ and $a$ such that for
every $s$ with $s(\theta)\ne0$ the definite trace of $f_s=(as)(\psi)*u^{*m}*e$ is nonzero. It remains to show that $f_s$ is good
(\cref{pf:ks1:def-spec}). Its support lies in $\mathcal A_m$, so it consists of regular elements satisfying $\mathrm R$. The cuspidality condition
$\mathrm C$ passes from $e$ to $u^{*m}*e$ and to $f_s$ (\lref{UpperTheory.Gal.Spc.cuspIdeal}). For a cuspidal function supported on
elements satisfying $\mathrm R$, the terms $\mathrm T_P$ of marked pairs $P$ with reducible characteristic polynomial vanish
(\lref{UpperTheory.Gal.Spc.vanPairs}) and the orbital integrals at elements of the definite groups which are not regular irreducible
vanish (\lref{UpperTheory.Gal.Spc.vanForms}); both rest on the vanishing of orbital integrals of cuspidal functions at non-elliptic
elements (\lref{UpperTheory.Gal.Spc.orbK_eq_zero_of_cuspSlices}). Hence $f_s$ is good and $\mathcal S$ holds with the scalar multiple
given by $a$.
\end{proof}

% ===========================================================================
\subsection{The holomorphic Lefschetz formula for finite \'etale correspondences}\label{sec:pf:geo:hlf}

Let $Z$ and $Y$ be schemes over $\C$. A \emph{complex point} of $Y$ is a morphism $y\colon\Spec\C\to Y$ over $\C$; we write $Y(\C)$ for
the set of complex points and $b(y)\in\C$ for the value of a regular function $b$ at $y$.

\begin{definition}[Tangent matrices]\label{pf:geo:istan}
Let $c_1,c_2\colon Y\to Z$ be morphisms over $\C$, $\delta\in\N$, $y\in Y(\C)$ with $c_1(y)=c_2(y)$, and $A\in M_\delta(\C)$. Then $A$ is a
\emph{tangent matrix of $(c_1,c_2)$ at $y$} if there are an affine open $U\subseteq Z$, an affine open $V\subseteq c_1^{-1}U\cap c_2^{-1}U$
containing $y$, functions $t_1,\dots,t_\delta\in\OO(U)$ and a $\C$-linear map $\varphi\colon\Omega_{\OO(V)/\C}\to\C^\delta$ such that
$\varphi(b\eta)=b(y)\varphi(\eta)$ for $b\in\OO(V)$, every $\eta$ with $\varphi(\eta)=0$ is a sum $\sum_ib_i\eta_i$ with $b_i(y)=0$,
\[
  \varphi\bigl(d(c_1^*t_l)\bigr)=e_l\quad\text{and}\quad\varphi\bigl(d(c_2^*t_j)\bigr)=(A_{jl})_{l}\qquad(1\le j,l\le\delta).
\]
\formal{GlobalTheory.W1.IsTan}
\end{definition}

Thus $\varphi$ identifies the cotangent space of $Y$ at $y$ with $\C^\delta$, the functions $c_1^*t_l$ are local coordinates at $y$, and
$d(c_2^*t_j)=\sum_lA_{jl}\,d(c_1^*t_l)$ at $y$: if $c_1$ is \'etale, $A=(\partial(t_j\circ f)/\partial t_l)$ is the Jacobian matrix, in the
coordinates $t$, of the local self-map $f=c_2\circ c_1^{-1}$ of $Z$ at its fixed point $c_1(y)$. This is the notion used in
\cref{pf:ks2:def:istan}.

\begin{theorem}[\cref{thm:galois:hlf}]\label{pf:geo:hlf}
Let $Z$ be a scheme over $\C$ which is smooth of relative dimension $\delta$ and projective, $Y$ a scheme over $\C$, and
$c_1,c_2\colon Y\to Z$ finite \'etale morphisms over $\C$. Let $\mathcal F\subseteq Y(\C)$ be a finite set such that $y\in\mathcal F$ if and only
if $c_1(y)=c_2(y)$, and for $y\in\mathcal F$ let $A_y$ be a tangent matrix of $(c_1,c_2)$ at $y$ with $\det(1-A_y)\ne0$. Then for every
$p\in\N$ the spaces $\geoHdg{Z}{p}{b}$ are finite-dimensional and
\begin{equation}\label{eq:pf:geo:hlf}
  \sum_{b=0}^{\delta}(-1)^b\,\tr\bigl((c_1)_*\circ(c_2)^*\mid\geoHdg{Z}{p}{b}\bigr)
  =\sum_{y\in\mathcal F}\frac{[X^p]\det(1+XA_y)}{\det(1-A_y)}
  =\sum_{y\in\mathcal F}\frac{\tr\bigl(\bigwedge^pA_y\bigr)}{\det(1-A_y)}.
\end{equation}
\formallabel{UpperTheory.Cb.lefschetz_fixedPoint_finite_etale_correspondence}\formallabel{GlobalTheory.T.lefschetz_fixedPoint_finite_etale_correspondence}\formallabel{UpperTheory.Gal.Int.Geo.weight_one_geometric_realization}
\end{theorem}

The proposition $\sfname{HLF}$ used in \cref{pf:geo:parts} and in \cref{pf:ks1:def-hlf,pf:ks2:main} is exactly this statement.
The second equality in \eqref{eq:pf:geo:hlf} is the identity $\det(1+XA)=\sum_p\tr(\bigwedge^pA)X^p$.

\subsubsection{Relation with the classical formula}

\begin{littheorem}[{Atiyah--Bott \cite[Thm.~4.12]{AtiyahBott1968}}]\label{pf:geo:lit:ab}
Let $X$ be a compact complex manifold, $\mathcal V$ a holomorphic vector bundle on $X$, $f\colon X\to X$ a holomorphic map with only simple fixed
points (i.e.\ $\det(1-df_P)\ne0$ at every fixed point $P$), and $\phi\colon f^*\mathcal V\to\mathcal V$ a holomorphic bundle homomorphism. Then
\[
  \sum_q(-1)^q\tr\bigl((f,\phi)^*\mid H^q(X,\mathcal V)\bigr)=\sum_{f(P)=P}\frac{\tr\phi_P}{\det_\C(1-df_P)}.
\]
\end{littheorem}

With $\mathcal V=\Omega^p$ and $\phi=\bigwedge^p({}^tdf)$ this gives
\[
  \sum_q(-1)^q\tr\bigl(f^*\mid H^q(X,\Omega^p)\bigr)=\sum_P\frac{\tr(\bigwedge^p{}^tdf_P)}{\det(1-df_P)};
\]
see also \cite{ToledoTong1978}.

\begin{remark}[What goes beyond the classical formula]\label{pf:geo:rem:hlflit}
(1) If $Y=Z$ and $c_1=\mathrm{id}$, \cref{pf:geo:hlf} is the case $\mathcal V=\Omega^p$ of \cref{pf:geo:lit:ab} for $f=c_2$, transported
to algebraic Hodge cohomology by GAGA \cite{Serre1956} (which identifies $H^q(Z,\Omega^p)$ with the cohomology of the analytic
sheaf, compatibly with pull-backs): the tangent matrix $A_y$ is then the Jacobian matrix of $f$ at $y$ in the coordinates $t$, the
bundle map $\bigwedge^p{}^tdf_y$ has the matrix $\bigwedge^p$ of its transpose, and traces of exterior powers and $\det(1-\cdot)$ are
invariant under transposition.

(2) For a correspondence with $c_1\ne\mathrm{id}$ the formula is not covered by \cref{pf:geo:lit:ab}, whose geometric endomorphisms are
defined by a single map \cite{AtiyahBott1967}. Atiyah and Bott remark \cite[p.~459]{AtiyahBott1968} that their formula for curves may be extended ``to any
self-correspondence of a curve $X$ in a plausible manner'', recovering the trace formulas of Eichler and Selberg; and Tu
\cite{Tu2023} formulates, for transversal holomorphic correspondences on compact complex manifolds, a conjectural holomorphic
Lefschetz formula, and lists as an open problem the computation of traces of Hecke correspondences on coherent cohomology. We
do not know a published proof of \cref{pf:geo:hlf} for general finite \'etale correspondences. The formal development proves it by
an algebraic argument which does not use \cref{pf:geo:lit:ab}. Its structure is described below; the central algebraic lemma,
the \emph{chart-lift theorem}, and the auxiliary lemmas on coordinates and \v{C}ech complexes are cited.
\end{remark}

\subsubsection{Reduction to integral schemes}

\begin{littheorem}[{Finiteness \cite[Thm.~3.2.1]{Grothendieck1961}, \cite[Thm.~III.5.2]{Hartshorne1977}}]\label{pf:geo:lit:fin}
If $Z$ is an integral scheme which is smooth of relative dimension $\delta$ and projective over $\C$, then every $\geoHdg{Z}{a}{b}$ is
finite-dimensional.
\formallabel{UpperTheory.finiteDimensional_hodgeCohomology_of_smooth_proj}
\end{littheorem}

Indeed $\Omega^a_{Z/\C}$ is coherent (locally free of finite rank), and the cohomology of a coherent sheaf on a proper scheme over a
field is finite-dimensional; the development proves this.

\begin{proposition}[Integral case]\label{pf:geo:hlfint}
\Cref{pf:geo:hlf} holds if $Z$ and $Y$ are integral.
\formallabel{UpperTheory.Cb.holomorphic_lefschetz_finiteEtale_integral}
\end{proposition}

\begin{proposition}[Reduction]\label{pf:geo:hlfred}
If \cref{pf:geo:hlf} holds whenever $Z$ and $Y$ are integral, then it holds in general.
\formallabel{UpperTheory.Core.block_fixed}\formallabel{UpperTheory.Cb.holomorphic_lefschetz_finiteEtale_correspondence}
\end{proposition}

\begin{proof}
\emph{Decomposition of $Z$.} $Z$ is projective, hence quasi-compact, and smooth over $\C$, hence its connected components are
integral; so $Z$ is the disjoint union of finitely many integral clopen subschemes $V_1,\dots,V_r$
(\lref{UpperTheory.Cb.exists_integral_clopen_decomposition_of_smooth}). Each $V_j$ is smooth of relative dimension $\delta$ and
projective (a clopen subscheme of a projective scheme is projective), so the spaces $\geoHdg{V_j}{p}{b}$ are finite-dimensional
(\cref{pf:geo:lit:fin}), and $\geoHdg{Z}{p}{b}=\bigoplus_j\geoHdg{V_j}{p}{b}$ by restriction, since sheaf cohomology commutes with finite
disjoint unions; this proves the finite-dimensionality. Put $W_{jk}=c_1^{-1}V_j\cap c_2^{-1}V_k$; these are clopen in $Y$, $Y$ is their
disjoint union, and the restrictions $c_1\colon W_{jk}\to V_j$ and $c_2\colon W_{jk}\to V_k$ are finite \'etale. Pull-back and trace are
compatible with restriction to clopen pieces, so under the decomposition of $\geoHdg{Z}{p}{b}$ the operator $(c_1)_*(c_2)^*$ is the block
matrix with $(j,k)$ entry $(c_1|_{W_{jk}})_*(c_2|_{W_{jk}})^*\colon\geoHdg{V_k}{p}{b}\to\geoHdg{V_j}{p}{b}$, and its trace is the sum of the
traces of the diagonal blocks. A fixed point $y$ lies in $W_{jj}$ for the $j$ with $c_1(y)\in V_j$, and a tangent matrix of $(c_1,c_2)$ at $y$
is a tangent matrix of the restrictions to the open subscheme $W_{jj}$ (\lref{UpperTheory.Cb.isTangentMatrix_iff_restrict_open}). So
\eqref{eq:pf:geo:hlf} for $(Z,c_1,c_2)$ is the sum over $j$ of \eqref{eq:pf:geo:hlf} for $(V_j,c_1|_{W_{jj}},c_2|_{W_{jj}})$, and we may assume $Z$
integral.

\emph{Decomposition of $Y$.} $Y$ is \'etale over $Z$, hence smooth, and finite over $Z$, hence quasi-compact; so $Y$ is the disjoint union of
finitely many integral clopen subschemes $Y_i$, and the restrictions of $c_1,c_2$ to $Y_i$ are finite \'etale. The trace along $c_1$ is the sum
of the traces along the $c_1|_{Y_i}$ (\lref{UpperTheory.Cb.hodgeTrace_eq_sum_clopen_decomposition}), pull-back along $c_2$ is compatible with
restriction, and each fixed point lies in exactly one $Y_i$, where the tangent matrices are unchanged. Hence \eqref{eq:pf:geo:hlf} for
$(Z,Y,c_1,c_2)$ is the sum over $i$ of the formulas for $(Z,Y_i,c_1|_{Y_i},c_2|_{Y_i})$, which hold by assumption.
\end{proof}

\begin{proof}[Proof of \cref{pf:geo:hlf}]
Combine \cref{pf:geo:hlfint,pf:geo:hlfred}.
\end{proof}

\subsubsection{The integral case}

We describe the proof of \cref{pf:geo:hlfint}; it is an assembly of the following lemmas of the formal development. Let $Z,Y$ be
integral, so $Y$ is smooth of relative dimension $\delta$ (being \'etale over $Z$), and $c_1,c_2$ induce $\C$-embeddings
$c_1^*,c_2^*\colon\C(Z)\to\C(Y)$ of function fields. For a family $x=(x_0,\dots,x_\delta)$ of elements of a subfield $F$ of $\C(Z)$ and a subset
$I\subseteq\{0,\dots,\delta\}$ let $R_I(Z)\subseteq\C(Z)$ be the ring of functions regular on $U_I=\bigcap_{i\in I}\{x_i\ne0\}$; similarly
$R_I(Y)\subseteq\C(Y)$ for the pulled-back sets $c_1^{-1}U_I$. We say that $x$ defines a \emph{finite cover} if the $U_i$ are affine opens covering
the scheme, with coordinate rings $R_{\{i\}}$, finite over the polynomial rings in the $x_m/x_i$.

\begin{enumerate}[label=(\alph*),nosep]
\item \emph{Coordinates} (\lref{GlobalTheory.T.exists_finiteCover_projSpace_reduced_fibre_origin}). There are $F$ and $x$ defining a
  finite cover of $Z$ (a finite morphism to $\mathbb P^\delta$), such that the scheme-theoretic fibre of $U_0$ over the origin, i.e.\ the
  spectrum of $R_{\{0\}}(Z)/(x_1/x_0,\dots,x_\delta/x_0)$, is reduced, and the points $c_1(y)$, $y\in\mathcal F$, lie in $U_0$ and in this fibre.
  The proof uses forms without common zeros and transversal at prescribed complex points
  (\lref{GlobalTheory.T.exists_forms_transversal_at_complexPoints}).
\item \emph{\v{C}ech description} (\lref{GlobalTheory.T.exists_hodge_cohomology_equiv_cech_differentials}). The Hodge cohomology
  $\geoHdg{Z}{a}{\bullet}$ is isomorphic to the cohomology of the \v{C}ech complex $I\mapsto\Omega^a_{R_I(Z)/\C}$ of the cover $(U_i)$, and similarly
  for $Y$ with the cover $(c_1^{-1}U_i)$, which again defines a finite cover with reduced fibre over the origin
  (\lref{UpperTheory.A.isFinCover_pull} and \lref{UpperTheory.etaleCorrespondence_fieldExtension_statement}); these isomorphisms are
  natural for pull-backs and for traces of finite \'etale morphisms.
\item \emph{Points} (\lref{UpperTheory.complexPoints_bijective_algHom_affineOpen} and \lref{UpperTheory.etaleCorrespondence_complexPoints_statement}).
  Complex points of $Y$ in $c_1^{-1}U_i$ correspond bijectively to $\C$-algebra maps $e\colon R_{\{i\}}(Y)\to\C$, and the fixed points correspond to
  the \emph{fixed characters} (those $e$ for which the two embeddings $c_1^*,c_2^*$ of $R_{\{i\}}(Z)$ give the same point). There are no fixed
  characters over $U_i$ for $i\ne0$, and the fixed characters over $U_0$ vanish on the $x_m/x_0$.
\item \emph{Local parameters} (\lref{UpperTheory.exists_regular_parameters_of_reduced_quotient} and \lref{UpperTheory.etaleCorrespondence_local_parameters_statement}).
  For each fixed character $e$ and each $m$ there are $c,g_1,\dots,g_\delta\in R_{\{0\}}(Y)$ with $e(c)\ne0$ and
  $c\cdot c_2^*(x_m/x_0)=\sum_lg_l\cdot c_1^*(x_l/x_0)$; the matrix $A'_e=(e(g_l)/e(c))_{m,l}$ is conjugate to the tangent matrix $A_y$ of the
  corresponding fixed point $y$. Conjugate matrices have the same $\det(1-A)$ and the same coefficients of $\det(1+XA)$.
\item \emph{Chart-lift theorem} (\lref{GlobalTheory.T.exists_chart_lift_of_finite_cover}). Let $F\subseteq K\subseteq K'$ be fields over $\C$,
  $x$ a family in $F$ defining finite covers for $K$ and $K'$, and $j_2\colon K\to K'$ a $\C$-embedding. Assume: the fibre over the origin of the
  cover of $K'$ is reduced; there are no fixed characters away from $U_0$; the fixed characters over $U_0$ vanish on the $x_m/x_0$; at every
  fixed character the ideal of the origin is generated, locally, by the $x_l/x_0$; and the matrices $A'_e$ of (d) satisfy
  $\det(1-A'_e)\ne0$, together with a compatibility between the charts of $K$ and $K'$. Let $u_q$ be endomorphisms of the \v{C}ech cohomology
  groups $H^q$ of $I\mapsto\Omega^a_{R_I(K)}$ which \emph{lift} to the correspondence: every cocycle $b$ has a lift $\tilde b$ to the \v{C}ech
  complex of the double cover of $K'$ indexed by pairs of index sets, which on the second family is the image of $b$ under $j_2$, and whose
  part on the first family has trace (from $K'$ to $K$) a cocycle representing $u_q[b]$. If the $H^q$ are finite-dimensional, then
  \[
    \sum_{q=0}^{\delta}(-1)^q\tr(u_q)=\sum_{e}\frac{[X^a]\det(1+XA'_e)}{\det(1-A'_e)},
  \]
  the sum running over the fixed characters $e$ in the fibre over the origin.
\end{enumerate}

\begin{proof}[Proof of \cref{pf:geo:hlfint}]
Finite-dimensionality is \cref{pf:geo:lit:fin}. Choose $F,x$ by (a) and transport $(c_1)_*(c_2)^*$ to the \v{C}ech cohomology by (b); by the
naturality in (b), the transported operators lift to the correspondence in the sense of (e)
(\lref{UpperTheory.etaleCorrespondence_family_statement}), with $K=\C(Z)$, $K'=\C(Y)$, the field embedding $c_1^*$ as the structure
map and $j_2=c_2^*$; the remaining hypotheses of (e) are provided by (a), (c) and (d)
(\lref{UpperTheory.etaleCorrespondence_sections_statement} and \lref{UpperTheory.etaleCorrespondence_differentials_statement}). By (c) the
fixed characters in the fibre over the origin are in bijection with $\mathcal F$, and by (d) the local terms of (e) are those of
\eqref{eq:pf:geo:hlf}. Since the alternating trace is invariant under the isomorphisms of (b), (e) gives \eqref{eq:pf:geo:hlf}.
\end{proof}

\begin{remark}\label{pf:geo:rem:hlfscope}
The finite set $\mathcal F$ in \cref{pf:geo:hlf} must be exactly the set of fixed points, and every fixed point must be transversal
($\det(1-A_y)\ne0$); the matrices $A_y$ are only constrained at the points of $\mathcal F$. Finite-dimensionality is part of the conclusion. In
the proof of \cref{thm:geoone} the formula is used only through $\sfname{KSHodge}$, for Hecke correspondences on Kuga--Sato varieties
composed with fibrewise multiplications (\cref{pf:ks1:def-hlf}).
\end{remark}

